\documentclass[11pt,a4paper]{amsart}
\usepackage{amsmath,amssymb,amsthm,mathtools}
\usepackage[utf8]{inputenc}
\usepackage[T1]{fontenc}
\usepackage{lmodern,microtype,graphicx}
\usepackage[hidelinks]{hyperref}
\usepackage{orcidlink,enumitem,booktabs,array}
\usepackage[margin=1in]{geometry}
\raggedbottom
\hypersetup{hypertexnames=false}
\AtBeginDocument{\begingroup\setbox0=\hbox{$\square\mathbb{R}$}\endgroup}
\theoremstyle{plain}
\newtheorem{theorem}{Theorem}[section]
\newtheorem{lemma}[theorem]{Lemma}
\newtheorem{proposition}[theorem]{Proposition}
\newtheorem{corollary}[theorem]{Corollary}
\theoremstyle{definition}
\newtheorem{definition}[theorem]{Definition}
\newtheorem{assumption}[theorem]{Assumption}
\theoremstyle{remark}
\newtheorem{remark}[theorem]{Remark}
\DeclareMathOperator{\Var}{Var}
\DeclareMathOperator{\Cov}{Cov}
\DeclareMathOperator*{\argmin}{arg\,min}
\newcommand{\R}{\mathbb{R}}
\newcommand{\N}{\mathbb{N}}
\newcommand{\EE}{\mathbb{E}}
\newcommand{\PP}{\mathbb{P}}
\newcommand{\QQ}{\mathbb{Q}}

\DeclareMathOperator{\tr}{tr}
\DeclareMathOperator{\chol}{chol}
\DeclareMathOperator{\Log}{Log}
\DeclareMathOperator{\diag}{diag}
\DeclareMathOperator{\diam}{diam}
\DeclareMathOperator{\TV}{TV}
\newcommand{\Hell}{\mathsf H}
\title[Leverage and structural information]{Leverage and Structural Information in Noisy Stochastic-Volatility Prices}
\author[J.~Vidal Llaurad\'o]{J.~Vidal Llaurad\'o\,\orcidlink{0009-0002-1918-5458}}
\subjclass[2020]{62M05, 62F12, 62F25, 62P05}
\keywords{Stochastic volatility, leverage, weak identification, noisy price observations, minimax estimation, honest confidence sets}
\hypersetup{
  pdftitle={Leverage and Structural Information in Noisy Stochastic-Volatility Prices},
  pdfauthor={J. Vidal Llaurad\'o},
  pdfkeywords={stochastic volatility, leverage, weak identification, noisy price observations, minimax estimation, honest confidence sets}
}
\date{}

\begin{document}
\begin{abstract}
Leverage can preserve information about a hidden volatility drift as
covariance shrinks, provided its response propagates through the latent
dynamics. For a bivariate Ornstein--Uhlenbeck model observed through noisy
log-price levels on a fixed grid, we prove exact identification at positive
covariance and adaptive squared drift risk of order
$\min\{1,[n(s^2+se^2)]^{-1}\}$, where $n$ counts observation blocks,
$s$ is the covariance scale and $e$ the leverage amplitude. Matching
lower bounds on specified drift neighborhoods use the full price law.
Dynamically exciting leverage thus permits learning where zero leverage
prevents uniform consistency. Nuisance parameters are unknown; confidence
coverage includes zero covariance and absent excitation. Precision requires
fixed excitation and residual covariance margins, fixed positive level
noise, a nonaliased drift branch and a sufficiently small, unevaluated
covariance ceiling.
\end{abstract}
\maketitle

\section{When leverage changes learnability}\label{sec:intro}

When volatility fluctuations are small, noisy prices may contain little
evidence about the dynamics generating them. The persistence and interaction
of hidden volatility factors can then be difficult to estimate, even in a
long sample. We use noisy
price levels, with hidden volatility, and ask when return--volatility
dependence makes those dynamics learnable. The answer depends on how the
signed leverage response propagates through the latent system.

The sample grows through $n$ observation blocks on a fixed grid. In the
bivariate Ornstein--Uhlenbeck model studied here, the state covariance is
$\Sigma=sV$, with $\tr V=2$; hence $s=\tr\Sigma/2$, while the state
diffusion amplitude is proportional to $\sqrt{s}$. Covariance responses
have size $s$, whereas signed leverage responses have size $\sqrt{s}\,e$,
where $e$ is the Frobenius leverage amplitude. If the dynamics propagate
that response across both state coordinates, leverage supplies information
about every entry of the drift matrix.

At positive covariance and with a fixed excitation margin at the truth,
an exact nonlinear moment inverse identifies the whole drift against
alternatives throughout the admitted parameter class, with all nuisance
parameters unknown. A single observable estimator adapts to the covariance
and leverage amplitudes. Its squared drift risk is bounded by a constant
times
\[
 \min\{1,[n(s^2+se^2)]^{-1}\}.
\]
Lower bounds for the full observed-price law match this rate on specified
drift neighborhoods. No estimator based on those prices improves its
worst-case order. For
$s_n=n^{-a}$, $1/2<a<1$, covariance-only information $ns_n^2$ tends to
zero, while $n(s_n^2+s_ne^2)$ diverges for fixed nonzero dynamically
exciting leverage. Under the neighborhood hypothesis of
Theorem~\ref{thm:main}, this separates worst-case failure of uniform drift
consistency from consistent estimation on the same noisy grid.

Confidence sets translate this information boundary into uncertainty about
a fitted drift. Our sets have honest coverage over the full admitted class,
including zero covariance and absent excitation. Their expected squared
contrast diameter has the same upper order at excited positive-scale
truths, with matching worst-case lower bounds on the specified
neighborhoods. These precision guarantees require fixed excitation and
residual covariance margins and a sufficiently small, unevaluated
covariance ceiling.

Signed leverage moments have long been used to study latent volatility.
The correlated exponential-OU response and its diffusion-amplitude order
occur in Masoliver and Perell\'o \cite{MasoliverPerello2005}.
Ahsan, Dufour and Rodriguez-Rondon \cite{ADR2025journal} study signed
moments and recovery of unknown leverage and latent dynamics in a discrete
scalar higher-order model with specified lag support and regular
fixed-parameter theory. Here an exact integrated-price calculation allows
unknown carry, covariance, loading and level noise; propagation determines
which matrix drift directions benefit from the response.

Inference under weak identification is also established. Cheng, Renault
and Sangrey \cite{ChengRenaultSangrey2024} study volatility risk prices
through leverage under reduced-form regularity assumptions supporting
uniform inference. Ahsan and Dufour \cite{AhsanDufour2024} develop
identification-robust stochastic-volatility tests using high-frequency
instruments and nuisance projection. Our confidence construction follows
these principles, with diameter governed by the information about the
hidden drift as covariance shrinks.

A moment inverse alone leaves open whether other features of the prices
could give a better rate. Abraham, Naulet and Gassiat
\cite{AbrahamNauletGassiat2022} connect weak moment geometry with
full-sequence information for finite-state hidden models. Gloter
\cite{Gloter2007} studies price-based drift inference, including noisy
prices on a shrinking grid over a growing horizon. Chong, Hoffmann, Liu,
Rosenbaum and Szymanski \cite{ChongEtAl2024} prove actual-price minimax
results for rough volatility in fixed-horizon infill with independent
drivers and without microstructure noise. Our lower bounds concern the
full noisy-price law on a fixed grid as covariance shrinks.
Structural-shock identification with measurement error provides broader
context in Gregoir and Meddahi \cite{GregoirMeddahi2025}.

\section{The noisy-price experiment}\label{sec:model}

Fix $N\ge20$, $\Delta>h>0$, $\delta=h/N$ and $r=\delta$. For $n\ge2$ the
observation times are all distinct elements of
\begin{equation}\label{eq:grid}
 \mathcal T_n=\{k\Delta+j\delta:0\le k\le n+1,\ 0\le j\le N\}.
\end{equation}
Both asset coordinates are observed:
\begin{equation}\label{eq:prices}
 P_i(t)=X_i(t)+\varepsilon_i(t),\qquad
 dX_i(t)=b_i\,dt+e^{Y_i(t)}dU_i(t),\quad i=1,2.
\end{equation}
The clean initial log prices are fixed known constants. There is exactly
one error per coordinate and physical time,
\[
 \varepsilon_i(t)\sim N(0,\nu_i),
\]
independent across coordinates and distinct times and independent of all
drivers. Reusing a time reuses its error. Here $\nu_i$ is a
variance. Known bounds are
\begin{equation}\label{eq:nuisance-bounds}
 |b_i|\le B_c,\qquad \mu_-\le\mu_i\le\mu_+,\qquad
 0<\nu_-\le\nu_i\le\nu_+<\infty.
\end{equation}
The carry $b$ is an unknown constant log carry; there is no
compensation term in \eqref{eq:prices}. The experiment is the sigma-field
of every noisy level in \eqref{eq:grid}, including the initial noisy level
and the observations adjoining the gaps. Grid spacings and noise bounds
remain fixed as $n$ grows and covariance shrinks.

The state is a two-sided stationary Gaussian OU process:
\begin{equation}\label{eq:state}
 dY=M(Y-\mu)dt+A\,dB,\qquad
 \Sigma=sV,\qquad A=\sqrt{s}A_0,\qquad
 A_0=\chol[-(MV+VM^\top)].
\end{equation}
Cholesky is lower triangular with positive diagonal. Its convention fixes
the loading coordinates. The stationary prehistory has the model-implied
law and is unobserved. At $s=0$ set $A=0$ and $Y=\mu$.
The state coordinates are fixed by the outputs $e^{Y_i}$; there is no
unknown output mixing.

The state driver $B$ is standard two-dimensional Brownian motion, and
\begin{equation}\label{eq:brownian}
 d\langle U\rangle=Q\,dt,\quad
 Q=\begin{pmatrix}1&q\\q&1\end{pmatrix},\qquad
 d\langle B,U\rangle=L\,dt.
\end{equation}
Thus the row of $L$ records a state-driver coordinate and the column a
return-driver coordinate. Fix $0<\kappa<1$, $q_{\max}<1$, and require
\begin{equation}\label{eq:residual}
 |q|\le q_{\max},\qquad Q-L^\top L\succeq\kappa I.
\end{equation}
The full $L$, including support, rank, orientation and amplitude, is
unknown. Taking traces gives $\|L\|_F^2\le2(1-\kappa)$. At $Q=I$,
$\|L\|_{\mathrm{op}}^2\le1-\kappa$, so rank-one amplitudes can reach
$\sqrt{1-\kappa}$, while full-rank Frobenius amplitudes can reach
$\sqrt{2(1-\kappa)}$.

Fix a compact class $\mathcal K$ of $(M,V)$ with
\begin{equation}\label{eq:class}
 V=V^\top,\quad \tr V=2,\quad v_-I\preceq V\preceq v_+I,
 \qquad -(MV+VM^\top)\succeq\omega_- I,
\end{equation}
where $v_-,\omega_->0$, on the prescribed branch
\begin{equation}\label{eq:strip}
 |\operatorname{Im}\lambda(M)|<\pi/\Delta.
\end{equation}
Indeed $D=V^{-1/2}MV^{1/2}$ satisfies
$D+D^\top\preceq-(\omega_-/v_+)I$. This proves strict stability and
uniform exponential decay. Compactness gives a strict nonaliasing margin.
Defective matrices are allowed. The Lyapunov equation in
\eqref{eq:state} gives stationary covariance $sV$ and confirms the
normalization $s=\tr\Sigma/2$.

The class $\Theta(s_0)$ consists of all parameters
$\theta=(M,V,s,L,q,\mu,b,\nu)$ satisfying these conditions and
$0\le s\le s_0$. Every listed parameter is unknown. Define the physical
state--return loading, transition and leverage amplitude by
\begin{equation}\label{eq:HZ}
 H=AL,\qquad Z=e^{M\Delta},\qquad e=\|L\|_F.
\end{equation}
For a fixed $\eta>0$, the performance subset $\Theta_\eta(s_0)$ is defined
by
\begin{equation}\label{eq:excitation}
 \sigma_{\min}([H,ZH])\ge\eta\|H\|_F.
\end{equation}
Here $\sigma_{\min}$ is the smaller row singular value, including zero.
At $H=0$ both sides vanish: zero leverage and the zero-scale boundary are
included. At $H\ne0$, the instantaneous and propagated response together
span both state coordinates. The condition depends on $H$, so an orthogonal
driver refactorization
$A\mapsto AO$, $L\mapsto O^\top L$ leaves it unchanged.

The estimator's candidates range over the full class $\Theta(s_0)$,
as does confidence coverage. The number $\eta$ indexes precision guarantees
only. Constants and the sufficiently small fixed covariance ceiling $s_0$
may depend on it, the grid and compact envelopes; no uniformity as $\eta$
or $\kappa$ tends to zero is asserted.

We use Frobenius norms for matrices and Euclidean norms for vectors. The
structural loss and a fixed nonzero contrast are
\begin{equation}\label{eq:loss}
 \|\Delta(\widehat M-M)\|_F^2,\qquad
 \psi_C(M)=\tr(C^\top\Delta M).
\end{equation}
For two parameter values put
$s_* =\max(s,s')$ and
$h_* =\max(\|\sqrt\Delta H\|_F,\|\sqrt\Delta H'\|_F)$, and define
\begin{align}\label{eq:comparison}
 d_\eta(\theta,\theta')={}&\|\Sigma-\Sigma'\|_F
 +\|\sqrt\Delta(H-H')\|_F
 +(s_*+h_*)\|\Delta(M-M')\|_F\notag\\
 &+\|\Delta(b-b')\|+\|\mu-\mu'\|+\|\nu-\nu'\|+|q-q'|.
\end{align}
No triangle property of this weighted comparison is needed. Rescaling time
to $\Delta=1$ replaces $(M,H,b)$ by
$(\Delta M,\sqrt\Delta H,\Delta b)$ and adds the known shift
$(\log\Delta)/2$ to each state mean. We restore these units in conclusions.

\section{Dynamic excitation and the information boundary}\label{sec:theorem}

Set $\rho_n(s,e)=\min\{1,[n(s^2+se^2)]^{-1}\}$ for $s>0$, and set
$\rho_n(0,e)=1$.

\begin{theorem}\label{thm:main}
For fixed $\mathcal K$, grid, nuisance bounds, residual margin and
excitation margin, there is a sufficiently small fixed $s_0>0$ with the
following properties.
\begin{enumerate}[label=\textup{(\Alph*)},leftmargin=*]
\item\label{part:id} If $\theta\in\Theta_\eta(s_0)$ and
$\theta'\in\Theta(s_0)$ have the same joint law of observed prices for
$n\ge2$, then at $s>0$ their $(M,\Sigma,L,q,\mu,b,\nu)$ agree, and hence
also $s,V$. At $s=0$ only $(\mu,b,q,\nu)$ remain in the law. No zero-scale
law equals a positive-scale law anywhere in $\Theta(s_0)$.
\item\label{part:upper} One measurable exact-moment estimator over the
full $\Theta(s_0)$, not supplied $s,L,\eta$ or any nuisance value,
satisfies
\begin{equation}\label{eq:joint-risk}
 \sup_{\theta\in\Theta_\eta(s_0)}
 \EE_\theta d_\eta(\widehat\theta_n,\theta)^2\le C/n.
\end{equation}
At each such truth of positive scale,
\begin{equation}\label{eq:drift-upper}
 \EE_\theta\|\Delta(\widehat M_n-M)\|_F^2
 \le C\min\{1,[n(s^2+se^2)]^{-1}\}.
\end{equation}
\item\label{part:lower} Suppose $\mathcal K$ contains $V=I$ and a full
four-dimensional drift neighborhood of
\begin{equation}\label{eq:coupled}
 M_0=\begin{pmatrix}-a_1&g\\g&-a_2\end{pmatrix},\qquad
 a_i>0,\quad0<g<\sqrt{a_1a_2}.
\end{equation}
The ball and $\eta$ are as in Proposition~\ref{prop:rank-one} below.
For its fixed-$(s,e)$ excited slices $\Theta_{\eta;s,e}$, for
$0<s\le s_0$ and $0\le e\le\sqrt{2(1-\kappa)}$,
\begin{equation}\label{eq:minimax}
 \inf_{\widetilde M}\sup_{\theta\in\Theta_{\eta;s,e}}
 \EE_\theta\|\Delta(\widetilde M-M)\|_F^2\asymp\rho_n(s,e).
\end{equation}
For every fixed nonzero $C$, scalar contrast risk has order
$\|C\|_F^2\rho_n(s,e)$. The infima range over all measurable estimators
using every noisy level. For $0\le e\le\sqrt{1-\kappa}$ the lower orders already hold
on the rank-one subfamily $L=eE_{12}$. Full-rank
$L=(e/\sqrt2)I$ supplies the entire feasible Frobenius range. These
subfamilies are not restrictions supplied to the estimator.
\item\label{part:confidence} For every fixed $\alpha\in(0,1)$ there are
observable confidence sets, with all nuisance parameters projected, whose coverage
is at least $1-\alpha$ uniformly over full $\Theta(s_0)$, including
nonexciting truths and $s=0$. At $\theta\in\Theta_\eta(s_0)$ their projected
contrast diameter satisfies
\begin{equation}\label{eq:diam-upper}
 \EE_\theta\diam(\mathcal I_{n,C})^2
 \le C_\alpha\|C\|_F^2\rho_n(s,e).
\end{equation}
Under the neighborhood assumption in \textup{(C)}, for fixed
$\alpha<1/4$ this is also the worst-case lower order for every honest
confidence set based on all prices.
\end{enumerate}
\end{theorem}

The theorem is an order comparison for squared loss, not an efficient
information or LAN assertion. The physical $H$ is in
\eqref{eq:joint-risk}; no uniform unweighted root-$n$ rate for $L$ is
asserted. The proofs of identification and upper inference are in
Sections~\ref{sec:moments}--\ref{sec:estimation}, and the all-price lower
proof is in Section~\ref{sec:information}.

Two examples clarify why propagation enters the rate. In the coupled
model below, rank-one leverage reaches both state coordinates and permits
recovery of all four drift entries. In the separate diagonal model of
Section~\ref{sec:contrasts}, with known loading support, rank-one leverage
accelerates inference on only one persistence parameter; the other retains
its covariance-only rate. The same instantaneous rank can therefore have
different consequences for structural learning.

\begin{proposition}\label{prop:rank-one}
The coupled example \eqref{eq:coupled}, with $V=I$ and $L=eE_{12}$,
satisfies \eqref{eq:excitation} on a full four-dimensional drift ball with
one positive normalized margin, uniformly for
$0<e\le\sqrt{1-\kappa}$. The same ball admits a positive uniform margin
for $L=(e/\sqrt2)I$, $0<e\le\sqrt{2(1-\kappa)}$, at $Q=I$.
\end{proposition}
\begin{proof}
By the two-dimensional Cayley--Hamilton calculus,
$e^{M\Delta}=uI+vM$. For distinct eigenvalues,
\[
 v=\frac{e^{\lambda_1\Delta}-e^{\lambda_2\Delta}}
          {\lambda_1-\lambda_2}.
\]
A vanishing numerator would require their difference to be a nonzero
$2\pi i/\Delta$ multiple, impossible on the strict strip
\eqref{eq:strip}. For a repeated eigenvalue and a nontrivial Jordan block,
the quotient limit is $v=\Delta e^{\lambda\Delta}\ne0$. A repeated,
diagonalizable two-dimensional matrix is scalar, in which case propagation
cannot improve loading rank. Thus for nonscalar $M$,
$[H,ZH]$ and $[H,MH]$ have the same row rank.

At \eqref{eq:coupled} the first column of $\chol(-2M_0)$ is
\[
 h_0=\begin{pmatrix}\sqrt{2a_1}\\-2g/\sqrt{2a_1}\end{pmatrix},
 \qquad
 \det[h_0,M_0h_0]=\frac{2g(a_1a_2-g^2)}{a_1}>0.
\]
Indeed $M_0h_0$ has entries
$-a_1\sqrt{2a_1}-2g^2/\sqrt{2a_1}$ and
$g\sqrt{2a_1}+2a_2g/\sqrt{2a_1}$.
Only the second column of $A_0L$ is nonzero, but it and its propagation
span the two rows. Continuity gives a closed drift ball, full in all four
entries, with negative definite symmetric part, a strict strip margin and
a positive normalized singular value. The factor $\sqrt{s}e$ cancels
from the normalized condition. For $L=(e/\sqrt2)I$, $A_0$ is uniformly
invertible on that ball, giving another positive normalized margin.
We choose $\eta$ below both margins and shrink the ball if necessary. The stated
amplitude ranges satisfy \eqref{eq:residual} with $Q=I$.
\end{proof}

This propagation algebra is classical reachability and shifted-response
realization; Oymak and Ozay \cite{OymakOzay2019} give an LTI identification
and conditioning antecedent with observed Gaussian inputs. Here the inputs
and the loading are not observed. At leading order $D_k=Z^{k-1}W$,
$W=F_1(M)H$. If $W$ has rank one, choose
a nonzero row $w^\top$ with $w^\top W=0$. The perturbation
$Z'=Z+uw^\top$ leaves $D_1,D_2$ unchanged when $W$ is fixed. If
$[W,ZW]$ has row rank two, $w^\top ZW\ne0$, and
$Z'^2W=Z^2W+uw^\top ZW$ changes the third response. A small admissible
perturbation can keep $W$ fixed by setting
$M'=\Delta^{-1}\Log Z'$ and $H'=F_1(M')^{-1}W$.
This is an ambiguity of the leading two-lag response, not an exact
observed-law equivalence. Section~\ref{sec:inverse} proves that the
three-lag mechanism survives the nonlinear moments.

\begin{corollary}\label{cor:separation}
Under the lower-neighborhood hypothesis of Theorem~\ref{thm:main}, along
$s_n=n^{-a}$ with $1/2<a<1$, zero leverage prevents uniform drift
consistency. The single estimator is consistent at truths with any fixed
nonzero leverage satisfying a fixed excitation margin on the same grid.
For $a=3/4$, its excited squared risk order is $n^{-1/4}$, with matching
worst-case lower order on the excited slices.
\end{corollary}
\begin{proof}
Here $ns_n^2=n^{1-2a}\to0$, so \eqref{eq:minimax} at $e=0$ stays bounded
below. For fixed $e>0$,
$n(s_n^2+s_ne^2)=n^{1-2a}+e^2n^{1-a}\to\infty$, and
\eqref{eq:drift-upper} tends to zero. At $a=3/4$ the dominating term is
$e^2n^{1/4}$. The sequences eventually lie below the fixed $s_0$.
The first statement is minimax worst-case impossibility; it is not a
pointwise lower claim at every zero-leverage drift.
\end{proof}

\begin{remark}\label{rem:finite-noise}
For fixed $n$, the joint law of all distinct noisy levels is the convolution
of the corresponding clean-price law with a Gaussian having strictly
positive diagonal covariance. The convolution density is everywhere
positive, since the Gaussian kernel is positive at every point. Thus any
two model laws are mutually absolutely continuous. If one estimator
recovered two different structural values almost surely under their
respective laws, the event of recovery of the first would have probability
one under both laws, contradicting recovery of the second. This argument
also applies to singular Brownian covariance, but is not a rate assertion
at perfect leverage. The nondegenerate asymptotic theorem above uses fixed
$\eta,\kappa$ and an unevaluated small covariance ceiling. Its observable
estimators and sets are existence constructions, not calibrated practical
confidence widths.
\end{remark}

\section{Exact moments of observed prices}\label{sec:moments}

The observable moments must distinguish the drift while allowing carry,
covariance, loading and level noise to vary. We derive their exact forms
before comparing their sensitivity to those parameters.

For $t\in\{r,h\}$ put
$R_i^t(k)=P_i(k\Delta+t)-P_i(k\Delta)$ and define
\[
 m_i=\EE R_i^h=b_ih,\qquad q_i^t=\EE(R_i^t)^2,\qquad
 w_i^t=q_i^t-(t/h)^2m_i^2,\qquad
 \lambda_i=e^{2\mu_i+2\Sigma_{ii}}.
\]
The raw second moment $q_i^t$ is distinct from the return correlation $q$.
It\^o isometry and independent endpoint errors give
\begin{equation}\label{eq:two-lengths}
 w_i^t=t\lambda_i+2\nu_i,\qquad
 \lambda_i=\frac{w_i^h-w_i^r}{h-r},\qquad
 \nu_i=\frac{h w_i^r-r w_i^h}{2(h-r)}.
\end{equation}
These identities do not require the return driver to be independent of
the state. For $k=1,2,3$ define
\begin{equation}\label{eq:cubic-def}
 T_{kij}=\EE[R_i^h(k)^2R_j^h(0)]-q_i^h m_j,\qquad
 D_{kij}=\frac{T_{kij}}{2h^2\lambda_i\sqrt{\lambda_j}}.
\end{equation}
For $k=1,2$ define
\begin{equation}\label{eq:fourth-def}
 G_{kij}=\frac{\EE[R_i^h(k)^2R_j^h(0)^2]-q_i^h q_j^h
                      -2m_jT_{kij}}{h^2\lambda_i\lambda_j}.
\end{equation}
An earlier linear martingale return can correlate with a later square,
requiring the unknown-carry correction $2m_jT_{kij}$. A later linear return
has zero mean against earlier variables.

The exact raw feature vector $F(\theta)\in\R^{27}$ consists of
\begin{equation}\label{eq:raw-features}
 \left(m_i,q_i^r,q_i^h; \EE[R_1^h(0)R_2^h(0)]; \EE[R_i^h(k)^2R_j^h(0)]_{k=1,2,3}; \EE[R_i^h(k)^2R_j^h(0)^2]_{k=1,2}\right),
\end{equation}
with $i,j=1,2$ in the mixed arrays. The counts are $6+1+12+8=27$.
Estimation and inference use this exact feature map.
The transformations \eqref{eq:two-lengths}--\eqref{eq:fourth-def} are
Lipschitz on its compact image: $\lambda_i$ is uniformly bounded above
and away from zero. Shared endpoint errors between $R_i^r$ and $R_i^h$
are retained in empirical dependence; no independence between those
returns is assumed.

Write, for $t\ge0$,
\begin{equation}\label{eq:cov-loading}
 C(t)=e^{Mt}\Sigma,\qquad K(t)=e^{Mt}H.
\end{equation}
In particular $K_{ij}(t-u)du=\Cov(Y_i(t),dU_j(u))$ for $t>u$.
The symbol $C(t)$ here is a covariance kernel, not the fixed contrast
matrix of \eqref{eq:loss}.

\begin{lemma}\label{lem:exact-moments}
For $z=k\Delta+t$, $\tau=z-u$,
\begin{equation}\label{eq:exact-cubic}
 D_{kij}=\frac1{h^2}\int_0^h\!\int_0^h
 K_{ij}(\tau)e^{2C_{ij}(\tau)-\Sigma_{jj}/2}\,du\,dt.
\end{equation}
Moreover $G_k=G_k^0+J_k$, where
\begin{equation}\label{eq:cov-moments}
 G_{kij}^0=\frac1{h^2}\int_0^h\!\int_0^h
                [e^{4C_{ij}(z-u)}-1]\,du\,dt
\end{equation}
and
\begin{align}\label{eq:J}
 J_{kij}=\frac4{h^2}\int_0^h\!\int_0^h\!\int_0^u
 &K_{ij}(z-u)[K_{jj}(u-v)+2K_{ij}(z-v)]\notag\\[-2pt]
 &\quad\cdot e^{-\Sigma_{jj}+C_{jj}(u-v)
                    +2C_{ij}(z-u)+2C_{ij}(z-v)}\,dv\,du\,dt.
\end{align}
The same-block cross moment satisfies
\begin{equation}\label{eq:cross-second}
 \EE[R_1^h(0)R_2^h(0)]-m_1m_2
 =h q\sqrt{\lambda_1\lambda_2}
       e^{\Sigma_{12}-(\Sigma_{11}+\Sigma_{22})/2}.
\end{equation}
\end{lemma}
\begin{proof}
We first establish the Gaussian exponential-tilting identity needed here.
Let $Z_0$ be a Gaussian linear functional of the two-sided drivers. For an
adapted integrand $f$ on a bounded interval, whenever the expressions are
integrable,
\begin{equation}\label{eq:tilt}
 \EE\left[e^{Z_0}\int f(u)dU_j(u)\right]
 =\int\frac{\Cov(dU_j(u),Z_0)}{du}\,
                  \EE[e^{Z_0}f(u)]\,du.
\end{equation}
After division by $\EE e^{Z_0}$, the left-hand side is an expectation under
the Gaussian tilt $\exp(Z_0-\EE Z_0-\Var(Z_0)/2)$. The driver
covariances are unchanged, and their means shift by their covariance
with $Z_0$. Centered increments remain independent of the preceding
sigma-field. For an adapted step process, the centered stochastic sum
therefore has zero expectation, leaving exactly the deterministic drift
shift in \eqref{eq:tilt}. Localization, H\"older and It\^o convergence
extend it to the integrands used below. The Gaussian state has uniformly
bounded fixed exponential moments, and the Burkholder--Davis--Gundy
(BDG) inequality supplies fixed return moments, so both sides converge
in the required $L^1$ spaces.

Let $N_j(u)=\int_0^u e^{Y_j(v)}dU_j(v)$ be the clean centered earlier
return. A squared later martingale increment, when multiplied by an
earlier measurable variable, can be replaced in expectation by its
quadratic-variation integral. This follows from It\^o's formula for its
square, since the remaining later stochastic integral has zero expectation.
For $z\ge\Delta>h$, \eqref{eq:tilt} with $Z_0=2Y_i(z)$ gives
\[
 \EE[e^{2Y_i(z)}N_j(h)]
 =2\int_0^h K_{ij}(z-u)\,
                   \EE[e^{Y_j(u)+2Y_i(z)}]\,du.
\]
The Gaussian expectation is
\[
 \EE[e^{Y_j(u)+2Y_i(z)}]
 =\lambda_i\sqrt{\lambda_j}
                    e^{2C_{ij}(z-u)-\Sigma_{jj}/2}.
\]
Integrating in $z=k\Delta+t$ proves \eqref{eq:exact-cubic} for the clean
centered returns.

For the fourth moment, It\^o's formula gives
\[
 N_j(h)^2=\int_0^h e^{2Y_j(u)}du
             +2\int_0^h N_j(u)e^{Y_j(u)}dU_j(u).
\]
Against $e^{2Y_i(z)}$ the first term has expectation
$\lambda_i\lambda_j e^{4C_{ij}(z-u)}$, producing
\eqref{eq:cov-moments} after subtracting the product of second moments.
Applying \eqref{eq:tilt} to the second term gives
\[
 4\int_0^h K_{ij}(z-u)
                 \EE[N_j(u)e^{Y_j(u)+2Y_i(z)}]du.
\]
A second application, now with $Z_0=Y_j(u)+2Y_i(z)$, gives
\[
 \EE[N_j(u)e^{Y_j(u)+2Y_i(z)}]
 =\int_0^u [K_{jj}(u-v)+2K_{ij}(z-v)]
            \EE[e^{Y_j(v)+Y_j(u)+2Y_i(z)}]dv.
\]
Dividing the last Gaussian exponential expectation by
$\lambda_i\lambda_j$ yields exactly
\[
 e^{-\Sigma_{jj}+C_{jj}(u-v)+2C_{ij}(z-u)+2C_{ij}(z-v)}.
\]
This proves the full triple integral \eqref{eq:J}.

The later and earlier observed block returns are, respectively,
$m_i+N_i+\xi_i$ and $m_j+N_j+\xi_j$.
The disjoint block errors $\xi_i,\xi_j$ are independent centered Gaussians
and independent of the clean paths. The later centered return has zero
mean against any earlier variable, and disjoint clean martingale returns
have zero cross second moment. Expanding the cubic thus leaves
$\EE[N_i^2N_j]$ after subtracting $q_i^hm_j$.
Expanding the raw fourth moment leaves
$\EE[N_i^2N_j^2]-\EE N_i^2\EE N_j^2$ and
$2m_j\EE[N_i^2N_j]$ after subtracting $q_i^h q_j^h$.
The latter is precisely the carry correction in \eqref{eq:fourth-def}.
Error variances in the two disjoint blocks cancel in these corrected
moments. This argument does not cancel shared errors within a block.

Finally the cross quadratic variation is
$q\int_0^h e^{Y_1(t)+Y_2(t)}dt$. Independent errors across the two
coordinates add no cross covariance. The Gaussian expectation
$e^{\mu_1+\mu_2+(\Sigma_{11}+\Sigma_{22})/2+\Sigma_{12}}$
proves \eqref{eq:cross-second}. Every integration is over a bounded
interval; H\"older, Gaussian exponential moments and BDG justify Fubini
and the localization limits above. In particular, we have not assumed a
fixed-mean conditional-normal return likelihood under leverage.
\end{proof}

These formulas exhibit $D_k=O(\sqrt{s}e)$, $G_k^0=O(s)$ and
$J_k=O(se^2)$. The leverage response order is classical
\cite{MasoliverPerello2005,ADR2025journal}; neither it nor a moment rank
calculation proves the all-price information order.

\section{A one-endpoint inverse}\label{sec:inverse}

We work in $\Delta=1$ units in this section. Write
$\delta_M=\|M-M'\|_F$, $\delta_\Sigma=\|\Sigma-\Sigma'\|_F$,
$\delta_H=\|H-H'\|_F$, and $\epsilon=\|F(\theta)-F(\theta')\|$.
Until excitation is explicitly invoked, both points can be anywhere in
the admitted compact class. By admissibility and bounded $A_0$,
\begin{equation}\label{eq:H-small}
 h_*^2\le C s_*,\qquad
 h_*=\max(\|H\|_F,\|H'\|_F),\quad s_*=\max(s,s').
\end{equation}
All constants below are independent of the true $s,e,n$.

\subsection{The covariance-pair inverse}

Define the matrix averaging kernels
\begin{equation}\label{eq:averaging}
 F_k(M)=\frac1{h^2}\int_0^h\!\int_0^h
                  e^{M(k\Delta+t-u)}du\,dt,
 \quad A_h(M)=\frac1h\int_0^h e^{Mt}dt.
\end{equation}
They obey $F_k=Z^{k-1}F_1$ and
$F_1=A_h(M)ZA_h(-M)$. Each factor is invertible on the drift class. The
scalar entire function $(e^{hz}-1)/(hz)$ has value one at zero and only
nonzero purely imaginary zeros; no Hurwitz eigenvalue, nor its negative,
can be such a zero. Functional calculus proves invertibility even at a
Jordan block, and compactness bounds $F_1,F_1^{-1}$ uniformly.

On the compact nonaliasing image,
\begin{equation}\label{eq:log-inverse}
 M=\Delta^{-1}\Log Z,\qquad \|M-M'\|_F\le C\|Z-Z'\|_F.
\end{equation}
Here $Z$ has spectrum off the nonpositive real axis by
\eqref{eq:strip}, and the principal logarithm of $e^{M\Delta}$ equals
$M\Delta$, also for defective $M$. The logarithm is analytic in a
neighborhood of each image point. A finite covering bounds its local
Fr\'echet derivatives and gives the Lipschitz inequality for close pairs.
For pairs outside those neighborhoods, continuity, injectivity and compact
separation give a positive minimum image distance, while the drift
diameter is bounded. This proves the global inequality; no spectral
diagonalization assumption is used.

\begin{lemma}\label{lem:cov-inverse}
For a sufficiently small fixed covariance ceiling,
\begin{equation}\label{eq:cov-inverse}
 \delta_\Sigma+s_*\delta_M
 \le C\|(G_1^0,G_2^0)-(G_1^{0\prime},G_2^{0\prime})\|.
\end{equation}
\end{lemma}
\begin{proof}
The linear terms in \eqref{eq:cov-moments} are
$B_1=4F_1\Sigma$ and $B_2=ZB_1$. We may label the endpoints so that
$s'=s_*>0$. Then $\sigma_{\min}(B_1')\ge c s_*$, and
\[
 (Z-Z')B_1'=(B_2-B_2')-Z(B_1-B_1')
\]
controls $s_*\|Z-Z'\|_F$, hence $s_*\delta_M$ by
\eqref{eq:log-inverse}, by the linear-pair difference. Also
\[
 B_1-B_1'=4F_1(\Sigma-\Sigma')+4(F_1-F_1')\Sigma'
\]
and bounded $F_1^{-1}$ control $\delta_\Sigma$ by that difference and
$s_*\delta_M$. If $s_*=0$ both terms to be bounded vanish.

For the exponential remainder $e^{4C}-1-4C$, finite-time differentiation
has covariance derivative $O(s_*)$ and drift derivative $O(s_*^2)$.
Thus its integrated pair difference is at most
$Cs_*[\delta_\Sigma+s_*\delta_M]$. These estimates can be checked along
bounded matrix interpolations: finite-time exponentials have bounded
Fr\'echet derivatives regardless of stability of the interpolation, and
$\|\Sigma\|\le Cs_*$ throughout. This does not require convexity of
$\mathcal K$. Applying the linear inverse to the nonlinear pair gives a
remainder that can be absorbed when $s_0$ satisfies $Cs_0\le1/2$.
\end{proof}

\subsection{Profiling the unknown physical loading}

\begin{lemma}\label{lem:profile}
Write $D_k=F_kH+E_k$, $k=1,2,3$. Uniformly,
\begin{align}\label{eq:E-bounds}
 \|E_k\|_F&\le Cs\|H\|_F,\notag\\
 \|E_k-E_k'\|_F&\le C[s_*\delta_H+h_*\delta_\Sigma
                                      +s_*h_*\delta_M].
\end{align}
For a sufficiently small fixed $s_0$, profiling $H$ through $D_1$ gives
\begin{align}\label{eq:H-profile}
 \delta_H&\le C[\epsilon+h_*\delta_M+h_*\delta_\Sigma],\notag\\
 \|J_k-J_k'\|_F&\le C[h_*\epsilon+h_*^2(\delta_M+\delta_\Sigma)].
\end{align}
In addition, with $a=\|H\|_F$ at the unprimed endpoint,
\begin{equation}\label{eq:amplitude}
 h_*\le C(a+\epsilon).
\end{equation}
No excitation of the primed endpoint is needed.
\end{lemma}
\begin{proof}
Subtracting one from the exponential in \eqref{eq:exact-cubic} leaves a
remainder of order $O(s)$. Its covariance derivative is bounded, and its
drift derivative is $O(s)$, on all the finite integration times involved.
The expression is linear in $H$, while the derivative of $e^{Mt}$ is
bounded at finite times. Taking differences gives \eqref{eq:E-bounds}.
Bounded interpolations as above justify these pairwise bounds on a
possibly nonconvex compact parameter class.

For each return column $j$ the exact first cubic is
$D_{1,:j}=\mathcal B_j(M,\Sigma)H_{:j}$, where
\begin{equation}\label{eq:profile-matrix}
 (\mathcal B_j)_{ip}=\frac1{h^2}\int_0^h\!\int_0^h
 [e^{M(\Delta+t-u)}]_{ip}
 e^{2C_{ij}(\Delta+t-u)-\Sigma_{jj}/2}\,du\,dt.
\end{equation}
It is $F_1+O(s)$. The preceding uniform inverse bound and a Neumann
argument give uniformly bounded inverses for small $s_0$. Their coefficient
functions are Lipschitz: locally this follows by differentiation of an
inverse, and globally by compact separation for distant pairs. Therefore
subtracting $H_{:j}=\mathcal B_j^{-1}D_{1,:j}$ gives the first bound in
\eqref{eq:H-profile}. Here differences of $D_1$ are bounded by $C\epsilon$
through the raw normalization, and $\|D_1\|_F\le C\|H\|_F$.

Substitution of the column profiles into \eqref{eq:J} gives
$\overline J_k(M,\Sigma,D_1)$ as a quadratic polynomial in the observable
$D_1$, with bounded Lipschitz coefficient functions. Its difference is
bounded by a linear term $Ch_*\epsilon$ and a coefficient difference
$Ch_*^2(\delta_M+\delta_\Sigma)$, proving the second bound.
Finally
\[
 \|H'\|_F\le C\|D_1'\|_F
 \le C(\|D_1\|_F+\epsilon)\le C(a+\epsilon),
\]
which proves \eqref{eq:amplitude}. In particular, an arbitrary alternative
loading is controlled by the observed first cubic, not assumed known or
excited.
\end{proof}

\begin{lemma}\label{lem:pre-cov}
For arbitrary admitted endpoints,
\begin{equation}\label{eq:pre-cov}
 \delta_\Sigma\le C[\epsilon+s_*\delta_M].
\end{equation}
\end{lemma}
\begin{proof}
From \eqref{eq:cov-moments},
$G_1^0=4F_1\Sigma+R_1^0$, with
\begin{equation}\label{eq:cov-remainder}
 \|R_1^0-R_1^{0\prime}\|_F
 \le C[s_*\delta_\Sigma+s_*^2\delta_M].
\end{equation}
We take the difference of $G_1=G_1^0+J_1$ between the endpoints and apply
the inverse of $F_1$.
The difference of $F_1$ contributes at most $Cs_*\delta_M$.
By \eqref{eq:H-profile}, \eqref{eq:H-small} and bounded $h_*$, the
correction contributes at most
$C\epsilon+Cs_*\delta_M+Cs_*\delta_\Sigma$.
Together with \eqref{eq:cov-remainder} this yields
$\delta_\Sigma\le C[\epsilon+s_*\delta_M]+Cs_0\delta_\Sigma$.
For fixed $Cs_0\le1/2$, the last term can be absorbed. The covariance
derivative of the profiled correction is small in $s$, so this step
requires no smallness assumption on $L$.
\end{proof}

\subsection{Lag three and excitation at only one endpoint}

\begin{lemma}\label{lem:dynamic-inverse}
For every $\theta\in\Theta_\eta(s_0)$ and every
$\theta'\in\Theta(s_0)$, after another fixed smallness choice for $s_0$,
\begin{equation}\label{eq:dynamic-inverse}
 h_*\delta_M\le C_\eta\epsilon.
\end{equation}
\end{lemma}
\begin{proof}
The dynamic lag residuals
$\mathcal R_k=E_{k+1}-ZE_k$, $k=1,2$, satisfy the exact identities
\begin{equation}\label{eq:lag-identities}
 D_2=ZD_1+\mathcal R_1,\qquad D_3=ZD_2+\mathcal R_2.
\end{equation}
By \eqref{eq:E-bounds},
\[
 \|\mathcal R_k-\mathcal R_k'\|_F
 \le C[s_*\delta_H+h_*\delta_\Sigma+s_*h_*\delta_M].
\]
The additional term $(Z-Z')E_k'$ has precisely the last order here.
The first bound of \eqref{eq:H-profile}, together with bounded $s_*$, gives
\begin{equation}\label{eq:lag-remainder}
 \|\mathcal R_k-\mathcal R_k'\|_F
 \le C[s_*\epsilon+h_*\delta_\Sigma+s_*h_*\delta_M].
\end{equation}

If $\epsilon\ge a=\|H\|_F$, then \eqref{eq:amplitude} gives
$h_*\le C\epsilon$. The compact drift diameter proves
\eqref{eq:dynamic-inverse} immediately. This case includes $a=0$,
so it includes zero leverage and zero scale without any assumption on the
alternative.

Suppose instead $\epsilon<a$. Then $a>0$ and $h_*\le Ca$.
At the unprimed endpoint only, excitation gives
\begin{align}\label{eq:observed-excitation}
 \sigma_{\min}([D_1,D_2])
 &\ge\sigma_{\min}(F_1[H,ZH])-Csa\notag\\
 &\ge c\eta a-Csa\ge c_\eta a,
\end{align}
after choosing $s_0\le c\eta/(2C)$.
We used that $F_1$ commutes with $Z$, hence
$[F_1H,F_2H]=F_1[H,ZH]$. The difference of \eqref{eq:lag-identities}
takes the form
\[
 (Z-Z')D_1=(D_2-D_2')-Z'(D_1-D_1')
                         -(\mathcal R_1-\mathcal R_1'),
\]
and similarly with $D_2,D_3$. Equations \eqref{eq:lag-remainder} and
$h_*\le Ca$ give
\begin{equation}\label{eq:transition-action}
 \|(Z-Z')[D_1,D_2]\|_F
 \le C[\epsilon+a\delta_\Sigma+s_*a\delta_M].
\end{equation}
For any matrix $T$, \eqref{eq:observed-excitation} implies
$c_\eta a\|T\|_F\le\|T[D_1,D_2]\|_F$. Together with
\eqref{eq:log-inverse}, \eqref{eq:transition-action} and
\eqref{eq:pre-cov} and bounded $a$, this gives
\[
 a\delta_M\le C_\eta\epsilon+C_\eta s_*a\delta_M.
\]
Choosing $C_\eta s_0\le1/2$ allows the last term to be absorbed. Since
$h_*\le Ca$ in this case, \eqref{eq:dynamic-inverse} follows. The argument
places no lower bound on the singular values of $[H',Z'H']$.
\end{proof}

\begin{proposition}\label{prop:full-inverse}
For every excited truth and every admitted alternative,
\begin{equation}\label{eq:full-inverse}
 d_\eta(\theta,\theta')\le C_\eta\|F(\theta)-F(\theta')\|.
\end{equation}
\end{proposition}
\begin{proof}
Lemma~\ref{lem:cov-inverse}, applied to $G^0=G-J$, together with the
profiled correction and Lemma~\ref{lem:dynamic-inverse}, gives
\[
 \delta_\Sigma+s_*\delta_M
 \le C_\eta\epsilon+C h_*^2\delta_\Sigma.
\]
Indeed the term $h_*^2\delta_M$ is bounded by $C h_*\epsilon$, and $h_*$
is uniformly bounded. By \eqref{eq:H-small}, a further fixed smallness
choice allows $Ch_*^2\delta_\Sigma$ to be absorbed. Then
\eqref{eq:H-profile} and \eqref{eq:dynamic-inverse} yield
\begin{equation}\label{eq:structural-inverse}
 \delta_\Sigma+\delta_H+(s_*+h_*)\delta_M\le C_\eta\epsilon.
\end{equation}
The regular nuisance parameters follow explicitly: $b_i=m_i/h$,
$\nu_i,\lambda_i$ come from \eqref{eq:two-lengths}, and
$\mu_i=(\log\lambda_i)/2-\Sigma_{ii}$. Finally, by
\eqref{eq:cross-second},
\[
 q=\frac{\EE[R_1^hR_2^h]-m_1m_2}{h\sqrt{\lambda_1\lambda_2}}
       e^{-\Sigma_{12}+(\Sigma_{11}+\Sigma_{22})/2}.
\]
These maps are uniformly Lipschitz on the compact moment image.
Restoring the $\Delta$ units specified after \eqref{eq:comparison} gives
\eqref{eq:full-inverse}.

If $s>0$, equality of features gives $M,\Sigma,H$ and all regular
nuisance parameters. Then $s=\tr\Sigma/2$, $V=\Sigma/s$, $A$ is determined by
\eqref{eq:state}, and $L=A^{-1}H$. At zero scale the price law is that of
a constant-carry Brownian log price with covariance
$\diag(e^\mu)Q\diag(e^\mu)$ and the given level noise, independent of
$M,V,L$. The regular moments identify its retained parameters.
A zero-scale endpoint is always in $\Theta_\eta(s_0)$. Taking it as the
unprimed endpoint in \eqref{eq:full-inverse} rules out equality with
any positive-scale law. Equality of the joint laws of observed prices implies
equality of the available features for $n\ge2$. Conversely equal retained
parameters construct equal laws. This proves Theorem~\ref{thm:main}(A).
It does not classify pairs of large nonexciting positive-scale models.
\end{proof}

\section{Adaptive estimation and honest sets}\label{sec:estimation}

The inverse bound converts moment error into drift error according to the
signal strength at the truth. Estimation over the full parameter class
then allows that strength to be unknown, while confidence coverage remains
valid even where the drift cannot be precisely estimated.

\begin{lemma}\label{lem:moment-variance}
The average of the 27 raw features over origins $k=0,\ldots,n-2$, with all
block indices in \eqref{eq:raw-features} shifted accordingly, is denoted by
$\widehat F_n$. A deterministic compact-class envelope $C_F$ satisfies
\begin{equation}\label{eq:moment-variance}
 \sup_{\theta\in\Theta(s_0)}
 \EE_\theta\|\widehat F_n-F(\theta)\|^2\le C_F/n.
\end{equation}
This includes nonexciting leverage and zero scale.
\end{lemma}
\begin{proof}
An origin uses at most four neighboring blocks through lag three, all
already in $\mathcal T_n$. There are $n-1\ge n/2$ origins. Every feature
is a polynomial of degree at most four in returns. For every fixed $p$,
BDG bounds a clean return by the $p/2$ moment of
$\int e^{2Y_i(t)}dt$ on a bounded interval. H\"older and Gaussian
exponential moments bound this uniformly; the bounded carry and Gaussian
errors preserve the bound. Thus all raw features have uniformly bounded
fixed $L^p$ norms.

For $s>0$, the normalized state
$s^{-1/2}V^{-1/2}(Y-\mu)$ has transition correlation matrix
$R_t=V^{-1/2}e^{Mt}V^{1/2}$, with
$\|R_t\|_{\mathrm{op}}\le Ce^{-ct}$ independent of $s$. For sufficiently
large $t$, its pair of standardized states has joint covariance
$\left(\begin{smallmatrix}I&R_t^\top\\R_t&I\end{smallmatrix}\right)$.
Against the product Gaussian law the trace term in the Gaussian KL formula
vanishes and the determinant term gives
\[
 \mathrm{KL}=-\tfrac12\log\det(I-R_tR_t^\top)\le Ce^{-2ct}.
\]
Pinsker therefore bounds the corresponding total variation by $Ce^{-ct}$.
For a stationary Markov state, conditioning past and future on their
boundary states and then applying data processing passes this bound to
absolute regularity of separated state segments. It also passes to
emission blocks: after the boundary, joint Brownian innovations are
independent of the past, even though state and return innovations within
the future are correlated by $L$. Separated observation-error sets are
independent. Hence the feature blocks have an exponential absolute-
regularity bound once their spans $3\Delta+h$ are separated.

One separated block can be coupled to an independent copy with
mismatch probability at most its absolute-regularity coefficient $\beta$.
For $p>2$, H\"older on the mismatch event gives
\[
 |\Cov(f,g)|\le C\beta^{1-2/p}\|f\|_p\|g\|_p.
\]
Indeed independence makes the covariance with the copy zero, and the
difference $g-g^*$ is supported on the mismatch event; the three
H\"older exponents are $p,p,(1-2/p)^{-1}$.
Summing the exponential covariance tail and controlling the finitely
many overlapping offsets by Cauchy--Schwarz yields a uniform absolute
covariance sum. At $s=0$ separated features depend on independent return
innovations and errors and have finite-range dependence. For an average
of $n-1$ origins, the double covariance sum is at most $C(n-1)$, proving
\eqref{eq:moment-variance} after summing over the 27 coordinates.
The fixed moment envelopes, overlap count and geometric tail give a
known deterministic upper envelope $C_F$; no long-run variance estimate
or volatility proxy is required.
\end{proof}

Let $(\theta_j)_{j\ge1}$ be a fixed enumeration of a countable
dense subset of full $\Theta(s_0)$. We define $\widehat\theta_n$ as the
first enumerated point satisfying
\begin{equation}\label{eq:MD}
 \|F(\theta_j)-\widehat F_n\|
 \le \inf_{l\ge1}\|F(\theta_l)-\widehat F_n\|+1/n.
\end{equation}
Continuity of the exact moment map makes the countable infimum equal to
the compact infimum. The first-index selection is measurable, since all
its inequalities involve measurable real functions and a countable
infimum. It exists by the definition of the infimum, and
\begin{equation}\label{eq:MD-error}
 \|F(\widehat\theta_n)-F(\theta)\|
 \le2\|\widehat F_n-F(\theta)\|+1/n.
\end{equation}
At an excited truth, Proposition~\ref{prop:full-inverse} applies to the selected
candidate regardless of its excitation. Lemma~\ref{lem:moment-variance} then proves
\eqref{eq:joint-risk}. Uniform bounds for $A_0,A_0^{-1}$ give
\begin{equation}\label{eq:H-scale}
 c\sqrt{s}e\le\|\sqrt\Delta H\|_F\le C\sqrt{s}e.
\end{equation}
The drift weight in \eqref{eq:comparison} is at least
$s+\|\sqrt\Delta H\|_F$. Its square is comparable, within fixed constants,
to $s^2+se^2$. Dividing the joint bound by the squared truth weight and
using the compact drift diameter gives \eqref{eq:drift-upper}. This proves
Theorem~\ref{thm:main}(B). The construction gives an observable, measurable
estimator based on the exact moment map over the full parameter class.
Its computational efficiency is not established.

For fixed $\alpha\in(0,1)$ define the full-class moment confidence set
\begin{equation}\label{eq:confidence-ball}
 \mathcal C_{n,\alpha}=\left\{\vartheta\in\Theta(s_0):
       \|\widehat F_n-F(\vartheta)\|
                        \le\sqrt{C_F/(\alpha n)}\right\}.
\end{equation}
Projection gives a drift set or the contrast set
$\mathcal I_{n,C}=\{\psi_C(M_\vartheta):\vartheta\in\mathcal C_{n,\alpha}\}$,
with zero diameter for an empty set. The continuous compact construction
makes these observable sets and their diameters measurable. Markov's
inequality applied to \eqref{eq:moment-variance} gives
\begin{equation}\label{eq:honesty}
 \inf_{\theta\in\Theta(s_0)}
 \PP_\theta\{\theta\in\mathcal C_{n,\alpha}\}\ge1-\alpha.
\end{equation}
This coverage includes every nonexciting truth and $s=0$. Projection
removes nuisance coordinates rather than treating them as known, as in
identification-robust inference \cite{AhsanDufour2024,ChengRenaultSangrey2024}.

At an excited truth, every candidate in \eqref{eq:confidence-ball} obeys,
even off the coverage event,
\[
 \|F(\vartheta)-F(\theta)\|
 \le Z_n:=\sqrt{C_F/(\alpha n)}+\|\widehat F_n-F(\theta)\|.
\]
The one-endpoint inverse bounds its dimensionless drift distance from the
truth by $C_\eta Z_n/(s+\|\sqrt\Delta H\|_F)$ at positive scale.
Two candidates thus have contrast distance at most twice this bound times
$\|C\|_F$. Since $\EE Z_n^2\le C_\alpha/n$, and the projection also has
the fixed compact diameter, \eqref{eq:diam-upper} follows. At zero scale
the bounded diameter still gives the stated bound; structural uncertainty
generally cannot shrink there. This proves the upper part of Theorem~\ref{thm:main}(D)
without a leverage or rank pretest.

\section{The information in all observed prices}\label{sec:information}

The moment estimator attains the proposed rate. To show that using other
features of the prices cannot improve its worst-case order, we bound the
separation between nearby drift models in the full observed-price law.
For dominated laws we use
$\Hell^2(P,P')=\int(\sqrt p-\sqrt{p'})^2$, with
$\TV(P,P')=\frac12\int|p-p'|\le\Hell(P,P')$.
The path parameter in this section is $\xi$, distinct from confidence
error probability $\alpha$.

\begin{lemma}\label{lem:all-price}
Fix $Q=I$, $b,\mu,\nu,s,e$ and $\overline L$ with
$\|\overline L\|_F=1$, and let $L=e\overline L$ satisfy
$e^2\|\overline L\|_{\mathrm{op}}^2\le1-\kappa$. Suppose
$(M_\xi,V_\xi)$ is a $C^2$ path on a compact interval, with bounded
first and second derivatives and uniform class and stability bounds on
a neighborhood. The complete noisy-level laws satisfy
\begin{equation}\label{eq:price-Hellinger}
 \Hell^2(P_\xi^{(n)},P_{\xi'}^{(n)})
 \le Cn(s^2+se^2)|\xi-\xi'|^2.
\end{equation}
The constant is uniform in $n,s,e$ under the fixed envelopes. No
excitation is needed. This is a structural-submodel bound, not a
full-nuisance equivalence of experiments.
\end{lemma}

The proof is given in the next four subsections. At $s=0$ the laws on this
submodel do not depend on $\xi$. We therefore carry out the differentiation
for $s>0$.

\subsection{A parameter-free common driver and hidden states}

Let $\tau^2=1-\kappa/2$ and $r_0=\kappa/2$. Independent standard
two-dimensional two-sided Brownian motions $W^0,B^0,W^1$ give the representation
\begin{equation}\label{eq:common-factor}
 U=\tau W^0+\sqrt{r_0}W^1,\qquad
 B=\tau^{-1}LW^0+\Gamma_e^{1/2}B^0,\qquad
 \Gamma_e=I-LL^\top/\tau^2.
\end{equation}
Since $\|L\|_{\mathrm{op}}^2\le1-\kappa$,
$\Gamma_e\succeq\kappa/(2-\kappa)I$.
Direct multiplication gives
$\Cov(dU)=I\,dt$, $\Cov(dB)=I\,dt$ and
$\Cov(dB,dU)=L\,dt$. The factorization preserves the law and the
physical Cholesky convention, with positive independent conditional
components for both state and returns.

Let $I_\ell$ denote the intervals between successive distinct grid times.
These include every gap. Their lengths lie in
\begin{equation}\label{eq:interval-bounds}
 d_-:=\min(\delta,\Delta-h)>0,\qquad
 d_+:=\max(\delta,\Delta-h)<\infty,
\end{equation}
and their number is $(n+2)N+n+1\le Cn$. We consider an enlarged experiment
in which $W^0$ is observed on the full stationary prehistory up to the final
observation time, together with every clean asset-by-interval increment.
Both state coordinates remain hidden. The law of this additional driver
is independent of $\xi$.
Starting from the fixed clean initial levels, cumulative summation of
these increments reconstructs every clean level. Adding one original
independent Gaussian error for each coordinate and distinct time
reconstructs all observations, including the initial noisy level and gap
observations. This parameter-independent kernel reproduces shared
endpoint-error dependence in any subsequently formed returns.

Given $W^0$, the conditional state mean is
\begin{equation}\label{eq:conditional-mean}
 m_\xi(t)=\mu+\frac{\sqrt{s}e}{\tau}
 \int_{-\infty}^t e^{M_\xi(t-u)}A_{0,\xi}\overline L\,dW^0(u),
\end{equation}
and its covariance is $sR_\xi$, where, for $t\ge u$,
\begin{align}\label{eq:conditional-cov}
 R_\xi(t,u)&=e^{M_\xi(t-u)}S_\xi,\notag\\
 S_\xi&=\int_0^\infty e^{M_\xi v}A_{0,\xi}\Gamma_e A_{0,\xi}^\top
                                      e^{M_\xi^\top v}dv.
\end{align}
For $t<u$, the kernel is obtained by transposition. Smoothness of Cholesky
on the uniformly positive diffusion cone and Duhamel's formula give
\begin{equation}\label{eq:R-derivative}
 \|\dot R_\xi(t,u)\|_F\le C(1+|t-u|)e^{-c|t-u|}.
\end{equation}
Indeed
$\partial_\xi e^{M_\xi v}
 =\int_0^v e^{M_\xi(v-w)}\dot M_\xi e^{M_\xi w}dw$
is bounded by $Cv e^{-cv}$ after reducing $c$ if necessary; differentiating
the integral for $S_\xi$ is justified by this integrable bound.
In general $\dot R_\xi(t,t)\ne0$, even when the unconditional covariance
is fixed: conditioning on the common driver changes the covariance.

Differentiating the adapted convolution \eqref{eq:conditional-mean}
gives, for every fixed $p$,
\begin{equation}\label{eq:m-derivative}
 \sup_t\|\dot m_\xi(t)\|_{L^p}\le C_p\sqrt{s}e.
\end{equation}
For example, with $x=m_\xi-\mu$,
\[
 dx=M_\xi x\,dt+\frac{\sqrt{s}e}{\tau}A_{0,\xi}\overline L\,dW^0,
 \quad
 d\dot m_\xi=(M_\xi\dot m_\xi+\dot M_\xi x)dt
       +\frac{\sqrt{s}e}{\tau}\dot A_{0,\xi}\overline L\,dW^0.
\]
This stable triangular linear system has stationary second moments of
order $se^2$, by the exponential convolution bounds. Gaussian moments
give all fixed $p$. The second parameter derivative is likewise bounded
in $L^2$, using the $C^2$ coefficients and differentiated kernels.
For each finite collection of time nodes, Fubini bounds the $H^2$ norm
in $\xi$ almost surely; the one-dimensional embedding $H^2\subset C^1$
gives a continuously differentiable version at these nodes.

Let $j$ now index an asset $i(j)$ and a complete-grid interval $I_j$;
there are $m\le Cn$ such emissions. Conditional on $Y,W^0$ they are
independent normals with
\begin{equation}\label{eq:conditional-emission}
 a_j=b_{i(j)}|I_j|+\tau\int_{I_j}e^{Y_{i(j)}(t)}dW^0_{i(j)}(t),
 \qquad d_j=r_0\int_{I_j}e^{2Y_{i(j)}(t)}dt.
\end{equation}
This follows from independence of $W^1$ from $(Y,W^0)$ and independent
$W^1$ coordinates on disjoint intervals. The random mean is a measurable
It\^o functional of the state and the additional observed driver, defined
by its left sums.

\subsection{Finite Gaussian differentiation and the full score}

For both assets, we use a common deterministic refinement of the full grid
and left sums with positive cell lengths. At a node $u$, let $w_u$ be its
cell length, $z_u$ the
appropriate $W^0$ increment, and $y_u$ the corresponding left state value.
Both coordinates at the union of physical nodes form one Gaussian vector.
Distinct asset coordinates at one time are distinct Gaussian nodes.
For emission $j$ the finite quadrature is
\begin{equation}\label{eq:quadrature}
 a_j=b_i|I_j|+\tau\sum_{u\in j}e^{y_u}z_u,
 \qquad d_j=r_0D_j,\qquad D_j=\sum_{u\in j}w_u e^{2y_u}.
\end{equation}
Its node derivatives are
\[
 \mathfrak m_u=\partial_{y_u}a_j=\tau e^{y_u}z_u,\qquad
 \mathfrak d_u=\partial_{y_u}d_j=2r_0w_u e^{2y_u}.
\]
Their own second derivatives are $\mathfrak m_u$ and $2\mathfrak d_u$.
The state-mean derivative $\dot m_u$ is different from $\mathfrak m_u$.

For the one-dimensional normal density $k_j(x;a_j,d_j)$, set
$Z_j=(x-a_j)/\sqrt{d_j}$, $M_j=k_{j,a}/k_j$ and $S_j=k_{j,d}/k_j$.
Direct differentiation of its exponential gives
\begin{equation}\label{eq:Hermite-identities}
 \begin{gathered}
 M_j=\frac{H_1(Z_j)}{\sqrt{d_j}},\quad
 S_j=\frac{H_2(Z_j)}{2d_j},\quad
 \frac{k_{j,aa}}{k_j}=\frac{H_2(Z_j)}{d_j},\quad
 \frac{k_{j,ad}}{k_j}=\frac{H_3(Z_j)}{2d_j^{3/2}},\quad
 \frac{k_{j,dd}}{k_j}=\frac{H_4(Z_j)}{4d_j^2}.
 \end{gathered}
\end{equation}
Here $H_1(z)=z$, $H_2(z)=z^2-1$, $H_3(z)=z^3-3z$ and
$H_4(z)=z^4-6z^2+3$ are probabilists' Hermite polynomials.
For a finite Gaussian vector of mean $m_\xi$ and covariance $sR_\xi$,
\begin{equation}\label{eq:Gaussian-diff}
 \partial_\xi\EE f(Y)
 =\sum_u\dot m_u\EE f_u(Y)
       +\frac{s}{2}\sum_{u,v}\dot R_{uv}\EE f_{uv}(Y).
\end{equation}
The identity follows by differentiating the Gaussian characteristic
function. In the density, the mean derivative corresponds to
$-\sum_u\dot m_u\partial_u\phi$, and the covariance derivative to
$(s/2)\sum_{u,v}\dot R_{uv}\partial_{uv}\phi$.
Integrating against a compactly supported smooth $f$ and integrating by
parts once or twice gives \eqref{eq:Gaussian-diff}. Equivalently these
identities follow directly by differentiating the finite Gaussian density.
The conditional covariance at the distinct nodes is positive definite:
$\Gamma_e$ and the state diffusion are positive, as are nonzero-time OU
innovation covariances. For our positive emission functions, we first
localize in the Gaussian nodes. The uniform moment and inverse-moment
bounds proved below justify removal of that localization. Thus the identity
applies to the finite Gaussian vector at each quadrature.

Let $k=\prod_j k_j$. The score coefficients are
\begin{equation}\label{eq:score-singletons}
 \begin{aligned}
 \mathsf A_j&=\sum_{u\in j}\dot m_u\mathfrak m_u,&
 \mathsf B_j&=\sum_{u\in j}\dot m_u\mathfrak d_u,\\
 \mathsf C_j&=\sum_{u\in j}\dot R_{uu}\mathfrak m_u,&
 \mathsf E_j&=\sum_{u\in j}\dot R_{uu}(2\mathfrak d_u),
 \end{aligned}
\end{equation}
and
\begin{equation}\label{eq:score-pairs}
 \begin{aligned}
 \mathsf C^{mm}_{jk}&=\sum_{u\in j,v\in k}\dot R_{uv}
                                        \mathfrak m_u\mathfrak m_v,&
 \mathsf C^{md}_{jk}&=\sum_{u\in j,v\in k}\dot R_{uv}
                                        \mathfrak m_u\mathfrak d_v,\\
 \mathsf C^{dm}_{jk}&=\sum_{u\in j,v\in k}\dot R_{uv}
                                        \mathfrak d_u\mathfrak m_v,&
 \mathsf C^{dd}_{jk}&=\sum_{u\in j,v\in k}\dot R_{uv}
                                        \mathfrak d_u\mathfrak d_v.
 \end{aligned}
\end{equation}
The conditional marginal density $p_\xi(x\mid W^0)$ has derivative
\begin{equation}\label{eq:full-score}
 \partial_\xi p_\xi=\EE[k\mathcal S\mid W^0],\qquad
 \mathcal S=\sum_j(\mathsf A_jM_j+\mathsf B_jS_j)
 +\frac{s}{2}\left\{\sum_j(\mathsf C_jM_j+\mathsf E_jS_j)
                                            +\mathcal Q\right\},
\end{equation}
where the complete quadratic part is
\begin{align}\label{eq:Q-score}
 \mathcal Q={}&\sum_j\left(
 \mathsf C^{mm}_{jj}\frac{k_{j,aa}}{k_j}
 +2\mathsf C^{md}_{jj}\frac{k_{j,ad}}{k_j}
 +\mathsf C^{dd}_{jj}\frac{k_{j,dd}}{k_j}\right)\notag\\
 &+2\sum_{j<k}\left(
 \mathsf C^{mm}_{jk}M_jM_k+\mathsf C^{md}_{jk}M_jS_k
 +\mathsf C^{dm}_{jk}S_jM_k+\mathsf C^{dd}_{jk}S_jS_k\right).
\end{align}
For a single node,
$\partial_u k=k(\mathfrak m_u M_j+\mathfrak d_u S_j)$.
For $u=v$, its second derivative has the additional terms
$\mathfrak m_u k_{j,a}+2\mathfrak d_u k_{j,d}$, giving
$\mathsf C_j,\mathsf E_j$. The product of first emission derivatives gives
the four pair types in \eqref{eq:score-pairs}; within one emission the two
mixed types coincide after summing, and between two emissions symmetry of
$\dot R$ doubles the $j<k$ sum. The mean part of
\eqref{eq:Gaussian-diff} gives $\mathsf A_j,\mathsf B_j$.
The formula therefore includes derivatives of both the random state means
and the conditional covariance diagonals.

Since the law of the additional observed driver is parameter-free,
conditional Cauchy--Schwarz (or Jensen for the posterior score) bounds joint marginal Fisher
information by
\begin{equation}\label{eq:score-Jensen}
 I_\xi\le\EE\int k\mathcal S^2\,dx.
\end{equation}
For fixed state and driver, the $Z_j$ are independent standard normals.
The generating identity
$\EE[e^{tZ-t^2/2}e^{uZ-u^2/2}]=e^{tu}$ gives
$\EE H_a(Z)H_b(Z)=a!\,\mathbf1_{a=b}$ by comparing coefficients.
Products on distinct coordinate supports are therefore orthogonal.
In \eqref{eq:full-score} the singleton $H_2$ terms from
$\mathsf B_j$, $\mathsf E_j$ and $\mathsf C^{mm}_{jj}$ overlap; the
$H_1$ terms from $\mathsf A_j$ and $\mathsf C_j$ also share a degree.
The overlapping terms are bounded using $(x+y)^2\le2x^2+2y^2$.
The remaining singleton terms of degrees three and four, and the terms
on distinct pair supports, are orthogonal.
Consequently
\begin{align}\label{eq:score-norm}
 \int k\mathcal S^2\,dx\le{}&
 C\sum_j\left(\frac{\mathsf A_j^2}{d_j}
                         +\frac{\mathsf B_j^2}{d_j^2}\right)
 +Cs^2\sum_j\left(\frac{\mathsf C_j^2}{d_j}
                         +\frac{\mathsf E_j^2}{d_j^2}\right)\notag\\
 &+Cs^2\sum_{j,k}\left[
 \frac{(\mathsf C^{mm}_{jk})^2}{d_jd_k}
 +\frac{(\mathsf C^{md}_{jk})^2}{d_jd_k^2}
 +\frac{(\mathsf C^{dm}_{jk})^2}{d_j^2d_k}
 +\frac{(\mathsf C^{dd}_{jk})^2}{d_j^2d_k^2}\right].
\end{align}
The constant is independent of the emission count: each singleton or pair
support has only a fixed number of terms. Grouping by these supports avoids
a crude bound of order $n^2$ from squaring the sum without using
orthogonality.

\subsection{Mesh-uniform chronological bounds}

Put
\[
 a_{jk}=\sup_{t\in I_j,u\in I_k}
                  |\dot R_{i(j),i(k)}(t,u)|.
\]
For every fixed $p\ge2$ the coefficients satisfy, uniformly in the common
positive left quadrature,
\begin{align}\label{eq:coefficient-bounds}
 \|\mathsf A_j\|_p+\|\mathsf B_j\|_p&\le C_p\sqrt{s}e,\notag\\
 \|\mathsf C_j\|_p+\|\mathsf E_j\|_p&\le C_pa_{jj},\notag\\
 \|\mathsf C^{mm}_{jk}\|_p+\|\mathsf C^{md}_{jk}\|_p
 +\|\mathsf C^{dm}_{jk}\|_p+\|\mathsf C^{dd}_{jk}\|_p
                                      &\le C_pa_{jk}.
\end{align}
Expectations in these norms include both the common driver and the residual
state. We prove each bound.

For $\mathsf A_j$, $\dot m_u e^{y_u}$ is measurable before its Brownian
increment. BDG for an adapted step integrand, followed by Minkowski and
H\"older, gives
\[
 \left\|\sum_{u\in j}\dot m_u e^{y_u}z_u\right\|_p
 \le C_p\left(\sum_{u\in j}w_u
                     \|\dot m_u e^{y_u}\|_p^2\right)^{1/2}
 \le C_p\sqrt{s}e.
\]
The last inequality uses \eqref{eq:m-derivative} at order $2p$, uniform
Gaussian exponential moments and $\sum w_u=|I_j|\le d_+$.
For $\mathsf B_j$ there is only an ordinary weighted sum:
\[
 \left\|\sum_{u\in j}w_u\dot m_u e^{2y_u}\right\|_p
 \le\sum_{u\in j}w_u\|\dot m_u e^{2y_u}\|_p
 \le C_p\sqrt{s}e.
\]
The factors $\tau$ and $r_0$ are fixed. The coefficient $\dot R_{uu}$ is
deterministic and bounded by $a_{jj}$; the same martingale argument bounds
$\mathsf C_j$, and the weighted-sum argument bounds $\mathsf E_j$.

For $\mathsf C^{mm}_{jk}$, ordering the nodes by physical left time divides
the sum into two strict-order triangles and equal physical times. In the
triangle where $u$ precedes $v$, the contribution, apart from fixed
constants, is the adapted outer stochastic sum
\[
 \sum_{v\in k} e^{y_v}z_v
      \left(\sum_{\substack{u\in j\\t_u<t_v}}
                    \dot R_{uv}e^{y_u}z_u\right).
\]
For each fixed later $v$, the inner kernel $\dot R_{uv}$ is deterministic
and bounded by $a_{jk}$. BDG at order $2p$ gives its norm at most
$C_pa_{jk}\sqrt{d_+}$. It is measurable before $z_v$.
Multiplication by $e^{y_v}$ is bounded in $L^p$ by H\"older and exponential
moments. A second BDG application to the outer sum gives $C_pa_{jk}$.
Exchanging the two nodes gives the same bound for the reversed triangle.
This remains valid across asset coordinates because the common
Brownian filtration makes the earlier inner sum adapted to the later
coordinate's increment.

At equal physical times both states are measurable before the shared cell
increments. For either the same or different Brownian coordinates,
$\|z_{i,u}z_{k,u}\|_p\le C_pw_u$. H\"older bounds the exponential left
factors, so the equal-time contribution is at most
$C_pa_{jk}\sum w_u\le C_pa_{jk}d_+$. This includes the same-coordinate
quadratic-variation terms and the cross-asset equal-time products.

For $\mathsf C^{md}_{jk}$, the stochastic node
$u$ can precede or follow the ordinary-time node $v$. In the first case,
the sum is
\[
 \sum_{v\in k}w_ve^{2y_v}
       \sum_{\substack{u\in j\\t_u<t_v}}
                      \dot R_{uv}e^{y_u}z_u.
\]
BDG bounds the inner sum at order $2p$ by $C_pa_{jk}$; H\"older and
Minkowski control the outer time sum. In the second case the inner time
sum $\sum_{t_v<t_u}\dot R_{uv}w_ve^{2y_v}$ is adapted at $u$ and has
$L^{2p}$ norm at most $C_pa_{jk}d_+$; the outer stochastic sum is bounded
by BDG. For equal physical times an additional time weight multiplies a
Brownian increment, and the bound is
$C_pa_{jk}\sum w_u^{3/2}\le C_pa_{jk}d_+^{3/2}$.
The transposed ordering proves the $dm$ bound. Finally two applications
of Minkowski and H\"older give the $dd$ bound from the two ordinary-time
sums and $\sum_{u\in j}w_u\sum_{v\in k}w_v\le d_+^2$.
This proves every assertion in \eqref{eq:coefficient-bounds}. The chronological
ordering accounts for the dependence of the state on $W^0$.

All inverse powers of the integrated variances are uniformly integrable.
Weighted Jensen gives
\begin{equation}\label{eq:inverse-moments}
 D_j\ge |I_j|\exp\left(\frac2{|I_j|}\sum_{u\in j}w_uy_u\right),
 \qquad \EE D_j^{-p}\le d_-^{-p}e^{C_p}.
\end{equation}
The weighted state average in this display is unconditionally Gaussian
with mean bounded by the state-mean envelope and variance at most
$Cs_0$: its weights are positive and sum to one, and all state
covariances are bounded by $Cs_0$. Taking its exponential moment proves
the negative-moment bound for every fixed $p$.

H\"older's inequality applied to \eqref{eq:score-norm}, with
\eqref{eq:coefficient-bounds} at sufficiently high fixed moment orders
and \eqref{eq:inverse-moments}, with $d_j=r_0D_j$, bounds the information
as follows. The state-mean terms contribute $Cmse^2$. Covariance singletons
and pairs contribute
$Cs^2\sum_{j,k}a_{jk}^2$. For example
$\EE(\mathsf C^{mm}_{jk})^2/(d_jd_k)\le Ca_{jk}^2$ follows by a fourth
coefficient moment and a higher inverse moment; the other denominator
powers are handled identically. Constants may depend on $r_0>0$.
By \eqref{eq:R-derivative},
\begin{equation}\label{eq:summability}
 \sum_{j,k}a_{jk}^2\le Cm\le Cn.
\end{equation}
Indeed $(1+x)e^{-cx}\le C_ce^{-cx/2}$, and physical intervals separated
by $l$ indices have distance at least $(l-1)d_-$ once $l\ge2$.
There are only two asset copies per interval, so for each $j$ the sum
in $k$ is bounded by a geometric series; adjacent and coincident
intervals contribute finitely many bounded terms. Equations
\eqref{eq:score-Jensen}--\eqref{eq:summability} prove the finite-quadrature
bound
\begin{equation}\label{eq:finite-information}
 I_\xi\le Cn(s^2+se^2),
\end{equation}
independent of refinement, emission count and the true small parameters.

\subsection{Exact increments and Hellinger contraction}

We justify both the parameter integration and the passage to the exact
stochastic-integral experiment. For a fixed finite quadrature, the
conditional node mean has the $C^1$ version established above, and the
conditional Gaussian covariance is positive definite. Differentiation
under the finite Gaussian integral, first with localization, gives
\eqref{eq:full-score}. On a compact parameter interval the higher moments
in \eqref{eq:coefficient-bounds} and \eqref{eq:inverse-moments} permit
removing localization by H\"older and dominated integral estimates.
The resulting positive marginal density is pointwise absolutely
continuous in $\xi$ for almost every driver and return vector. The
integrable bounds on the localized derivative allow the fundamental
theorem of calculus to pass through the finite mixture by Fubini.
The driver law itself does not change with $\xi$.

Relative to that law times Lebesgue measure for returns,
$\partial_\xi\sqrt{p_\xi}=(\partial_\xi p_\xi)/(2\sqrt{p_\xi})$ has
squared $L^2$ norm $I_\xi/4$. The Jensen and moment bounds control this
norm uniformly on the parameter interval. Thus pointwise absolute
continuity, Fubini and Minkowski justify integrating the square-root
derivative in the joint $L^2$ space, giving
\begin{equation}\label{eq:Hellinger-integral}
 \Hell(p_\xi,p_{\xi'})
 \le\frac12\int_{\xi}^{\xi'}\sqrt{I_u}\,|du|
 \le C\sqrt{n(s^2+se^2)}|\xi-\xi'|.
\end{equation}
This derives the Hellinger bound from the finite-dimensional score.

For each fixed $n,\xi$, the positive left sums for $D_j$ converge almost
surely to the strictly positive integrals in
\eqref{eq:conditional-emission}, since the state paths are continuous.
The left mean sums converge in $L^2$ by It\^o isometry: the integrands
converge pointwise by continuity, and Gaussian exponential moments
uniformly dominate their squared differences on the finite horizon.
Normal kernels with means converging in probability and variances
converging to positive values converge in total variation in probability.
The deterministic kernel statement follows from pointwise density
convergence and Scheff\'e's lemma. For the finite product of emissions the
same conclusion holds. Its total variation is bounded, so its expectation
tends to zero. Coupling the conditional kernels with the same $(Y,W^0)$
therefore gives total-variation convergence of the joint law of the
emissions and the additional observed driver path.

Since $\Hell^2\le2\TV$, the triangle inequality and this convergence pass
the mesh-uniform bound \eqref{eq:Hellinger-integral} to exact richer
observations. Hellinger contracts under a parameter-independent kernel:
conditional Cauchy--Schwarz increases the affinity
$\int\sqrt{pp'}$ after a kernel, and
$\Hell^2=2-2\int\sqrt{pp'}$. Applying the reconstruction and noise kernel after
\eqref{eq:common-factor} gives \eqref{eq:price-Hellinger} for every
original noisy level. This completes the proof of
Lemma~\ref{lem:all-price}.

\subsection{Two-point risk and honest-diameter lower bounds}

We work on the drift ball in Proposition~\ref{prop:rank-one}, with $V=I$, $q=0$,
and fixed admissible $b,\mu,\nu,s$ and signed $L$.
For $0\le e\le\sqrt{1-\kappa}$ use $L=eE_{12}$; for the whole range
$0\le e\le\sqrt{2(1-\kappa)}$ use $L=(e/\sqrt2)I$.
Fix any Frobenius-unit direction $D$ and a short line
$M_\xi=M_0+\xi D$ inside the full four-dimensional ball. Its symmetric
part remains uniformly negative, so
$A_{0,\xi}=\chol[-(M_\xi+M_\xi^\top)]$ is $C^2$ with bounded
derivatives. Class, strip and excitation conditions hold uniformly along
the line. At $e=0$ the excitation condition holds by its zero convention.
The nuisance restrictions are used solely to construct the lower submodel,
not to inform the estimator.

In $\Delta=1$ units choose two interior points separated by
\begin{equation}\label{eq:two-point-separation}
 d=c\min\{1,[n(s^2+se^2)]^{-1/2}\},
\end{equation}
with fixed $c$ small enough for both to stay in the ball and for
Lemma~\ref{lem:all-price} to give $\TV\le1/2$.
Any estimate of two values a distance $d$ apart defines a test by selecting
the nearer value. On a classification error its squared loss is at least
$d^2/4$. The sum of error probabilities of any test is at least
$1-\TV$, so the maximum squared risk is at least
$d^2(1-\TV)/8$. This classical two-point argument applies to every
measurable function of all observed levels and proves the matrix lower
bound in \eqref{eq:minimax}. Restoring fixed $\Delta$ units changes only
constants. For a nonzero contrast choose $D=C/\|C\|_F$ in dimensionless
drift units. Its contrast separation is $\|C\|_F d$, proving the scalar
order including either off-diagonal entry. The upper order follows from
\eqref{eq:drift-upper} and Cauchy--Schwarz. This proves
Theorem~\ref{thm:main}(C), including the rank-one lower subfamily.

Let a confidence set be honest at these two points, and let their contrast
values be $\psi_0,\psi_1$. Under the first law,
\begin{equation}\label{eq:two-point-diameter}
 \PP_0\{\psi_0,\psi_1\text{ both in the set}\}
 \ge1-2\alpha-\TV(P_0,P_1)\ge1/2-2\alpha.
\end{equation}
Coverage of $\psi_1$ transfers from $P_1$ to $P_0$ with loss at most
$\TV$; intersecting with coverage of $\psi_0$ proves the inequality.
For fixed $\alpha<1/4$, the expected squared diameter is consequently
at least $(1/2-2\alpha)|\psi_1-\psi_0|^2$. This proves the matching
worst-case honest-diameter lower bound in Theorem~\ref{thm:main}(D).
At $s=0$, two distinct drifts with identical regular nuisance parameters have
the same law. The same event has probability at least $1-2\alpha$,
so honest structural uncertainty cannot uniformly shrink. Together with
the previously proved upper statements, this finishes
Theorem~\ref{thm:main}.

\section{Absent and nonexciting leverage}\label{sec:contrasts}

The same price experiment gives two contrasts in how leverage affects
learnability. With small unknown leverage, including zero, the
covariance-only rate remains attainable and has a matching worst-case bound
on a specified subfamily. In a diagonal state model with known
cross-loading support, the leverage term improves the rate for one
persistence while the other retains the covariance-only rate.

\subsection{Small unknown full leverage: a worst-case covariance boundary}

Let $\mathcal G(s_0,\ell_0)$ contain all parameters of the original compact
model with $\|L\|_F\le\ell_0$ and a fixed positive residual margin, allowing
$q=0,L=0$. Its loading support and orientation are unknown, and it has no
excitation condition.

\begin{proposition}\label{prop:small-leverage}
For sufficiently small fixed $\ell_0,s_0>0$, all positive-scale parameters
in $\mathcal G(s_0,\ell_0)$ are identified, zero and positive scales are
separated, and a single full-class exact-moment estimator has squared
dimensionless drift risk at most
$C\min\{1,(ns^2)^{-1}\}$. If the drift/shape class contains
\begin{equation}\label{eq:small-lower-subfamily}
 M_\gamma=\begin{pmatrix}-a_1&\gamma\\\gamma&-a_2\end{pmatrix},\quad
 V=I,\quad a_1\ne a_2>0,\quad
 \gamma\in J\Subset(0,\sqrt{a_1a_2}),
\end{equation}
where $J$ has positive length, this is the matching worst-case order at
fixed $s$, also for the off-diagonal contrast on that subfamily.
\end{proposition}
\begin{proof}
We again work in $\Delta=1$ units. The first-cubic profile of
Lemma~\ref{lem:profile} remains valid without excitation. On this class
$h_*\le C\sqrt{s_*}\ell_0$. Lemma~\ref{lem:cov-inverse}, applied to
$G^0=G-J$, and \eqref{eq:H-profile} give
\[
 \rho_0:=\delta_\Sigma+s_*\delta_M
 \le C\epsilon+C h_*\epsilon
                  +Ch_*^2(\delta_M+\delta_\Sigma)
 \le C\epsilon+C\ell_0^2(1+s_0)\rho_0.
\]
After fixing $s_0$ for the cubic profile and covariance inverse, we choose
$\ell_0$ so that the last coefficient is at most $1/2$. Absorbing the last
term gives $\rho_0\le C\epsilon$. The same regular formulas as before
recover $b,\nu,\lambda,\mu,q$. At positive scale, equality gives $M,\Sigma$, then
$H$ from the first cubic and $L=A^{-1}H$. At zero scale only the regular
parameters remain, and the bound on $\delta_\Sigma$ separates it from
positive scale.

The full-class raw variance lemma and dense minimum-distance construction,
now over $\mathcal G$, give a uniform $C/n$ squared bound for $\rho_0$
plus regular nuisance errors. The truth drift weight is at least $s$;
division and compact capping give the upper rate. The extra cubic lag in
our 27 features is harmless; the inverse only needs the covariance pair
and first cubic.

For the lower bound, we use \eqref{eq:small-lower-subfamily}, $L=0,Q=I$ and
fixed admissible nuisance parameters and $s$. This symmetric drift submodel
contains an interval of positive length. Lemma~\ref{lem:all-price}
at $e=0$ gives $\Hell^2\le Cns^2|\gamma-\gamma'|^2$;
its proof here is covariance-only since the conditional state mean is
fixed. Choosing an interior pair of separation
$c\min\{1,(\sqrt n s)^{-1}\}$ and applying the all-price two-point test
proves the lower order, also for $M_{12}=\gamma$.
\end{proof}

This worst-case result uses the admitted zero-leverage subfamily. It is
not a pointwise impossibility claim at every small nonzero loading.
Dynamic excitation in Theorem~\ref{thm:main} gives another inverse route,
without a small-leverage condition on its performance subset.

\subsection{A known-support nonexciting family}

We next consider the known restrictions
\begin{equation}\label{eq:diagonal-family}
 \begin{gathered}
 M=\diag(-a_1,-a_2),\qquad \Sigma=\diag(v_1,v_2),\qquad
 v_i=sV_i,\quad V_1+V_2=2,\\
 A=\diag(\sqrt{2a_1v_1},\sqrt{2a_2v_2}),\qquad L=\ell E_{12}.
 \end{gathered}
\end{equation}
The $a_i$ range independently over fixed positive compact intervals of
positive length. The $V_i$ range over compact positive shapes with
$(1,1)$ admitted. Signed $\ell$, scale, shape, $q,\mu,b,\nu$ are all
unknown, with the previous nuisance bounds and
\begin{equation}\label{eq:diagonal-margin}
 q^2+\ell^2\le1-\kappa.
\end{equation}
This is a separate compact class $\mathcal D(s_0)$: its diagonal state
structure and leverage support are known, not an arbitrary unknown
rank-one restriction. The determinant and trace of
$Q-L^\top L=\left(\begin{smallmatrix}1&q\\q&1-\ell^2\end{smallmatrix}\right)$
show that its least eigenvalue is at least $\kappa/2$, so all required
uniform residual bounds hold with that fixed margin. At $\ell\ne0$,
$H$ and $ZH$ have only their first rows nonzero. Thus their concatenation
has row rank one and does not satisfy \eqref{eq:excitation}.

\begin{theorem}\label{thm:diagonal}
For a sufficiently small fixed $s_0$, all positive-scale parameters in
$\mathcal D(s_0)$ are exactly identified. One exact-moment estimator over
the whole class, with all parameters unknown, attains simultaneously
\begin{equation}\label{eq:directional-rates}
 \begin{gathered}
 \EE(\widehat a_1-a_1)^2\le Cr_1,\qquad
 \EE(\widehat a_2-a_2)^2\le Cr_2,\\
 r_1=\min\{1,[n(s^2+s\ell^2)]^{-1}\},\quad
 r_2=\min\{1,(ns^2)^{-1}\}.
 \end{gathered}
\end{equation}
On each slice $s>0$, $|\ell|=e\le\sqrt{1-\kappa}$, these are matching
minimax orders for every estimator using all noisy levels. For each fixed
$c=(c_1,c_2)$ the contrast $c_1a_1+c_2a_2$ has order
$c_1^2r_1+c_2^2r_2$. Full-class honest projected confidence sets have
this expected squared diameter upper order for every fixed
$\alpha\in(0,1)$, and matching worst-case lower order for fixed
$\alpha<1/4$.
\end{theorem}

We prove the upper inverse first and then the two separate all-price
obstructions; the latter cannot be inferred from the moment derivatives.

\subsection{Exact directional inverse and upper bounds}

In \eqref{eq:diagonal-family},
$K_{12}(t)=\sqrt{2a_1v_1}\ell e^{-a_1t}$, $C_{12}(t)=0$, and
$K_{11}=K_{22}=0$. Lemma~\ref{lem:exact-moments} gives, exactly,
\begin{equation}\label{eq:diagonal-cubic}
 T_{k,12}=2h^2\lambda_1\sqrt{\lambda_2}e^{-v_2/2}
          \sqrt{2a_1v_1}\ell f_h(a_1)e^{-a_1k\Delta},
 \qquad f_h(a)=\frac{2(\cosh(ah)-1)}{a^2h^2}>0.
\end{equation}
The factor is $f_h(a)=A_h(-a)A_h(a)$ by direct integration.
Both own-coordinate corrections $J_{k,ii}$ vanish, since their first
loading factor $K_{ii}$ is zero. The own fourth-moment pairs are therefore
exactly their scalar covariance-only maps, for every $q,\ell$.

The scalar proof of Lemma~\ref{lem:cov-inverse} gives, for any two
candidates in this entire restricted class,
\begin{equation}\label{eq:scalar-inverse}
 |v_i-v_i'|+\max(v_i,v_i')|a_i-a_i'|
          \le C\|F(\theta)-F(\theta')\|,
 \qquad i=1,2.
\end{equation}
Specifically the scalar linear pair is
$B_1=4v_i f_h(a_i)e^{-a_i\Delta}$, $B_2=e^{-a_i\Delta}B_1$;
the larger variance endpoint has $|B_1'|\ge c\max(v_i,v_i')$.
The scalar logarithm of $B_2/B_1$ controls persistence on the positive
compact interval. Covariance and drift remainder derivatives are
$O(v_*)$ and $O(v_*^2)$, as before, and are absorbed for small $s_0$.

Put $z(a)=e^{-a\Delta}$. Equation \eqref{eq:diagonal-cubic} gives
$T_{2,12}=z(a_1)T_{1,12}$ and, for any alternative,
\[
 |T_{1,12}|\,|z(a_1)-z(a_1')|
 \le|T_{2,12}-T_{2,12}'|
          +|z(a_1')|\,|T_{1,12}-T_{1,12}'|.
\]
On the compact interval $|z(a)-z(a')|\ge c|a-a'|$, and the exact
amplitude obeys $|T_{1,12}|\ge c\sqrt{v_1}|\ell|$. The centered
transformations are Lipschitz in the raw features. Thus, without knowing
magnitude or sign at either endpoint,
\begin{equation}\label{eq:scalar-leverage-inverse}
 \sqrt{v_1}|\ell|\,|a_1-a_1'|\le C\|F(\theta)-F(\theta')\|.
\end{equation}
Together with \eqref{eq:scalar-inverse}, this gives the truth-side bounds
\begin{equation}\label{eq:directional-inverse}
 \sqrt{v_1^2+v_1\ell^2}\,|a_1-a_1'|\le C\|F-F'\|,
 \qquad v_2|a_2-a_2'|\le C\|F-F'\|.
\end{equation}
They include comparison with zero-scale alternatives. Since $v_i$ is
comparable to $s$, Lemma~\ref{lem:moment-variance} and the measurable
full-class minimum-distance construction prove both upper rates in
\eqref{eq:directional-rates}. In this statement full class means the
entire class $\mathcal D(s_0)$, not the unrestricted class with unknown
loading support. One estimator attains both rates simultaneously.

At positive scale \eqref{eq:scalar-inverse} identifies $a_i,v_i$,
including at $\ell=0$. Equation \eqref{eq:diagonal-cubic} then identifies
the signed $\ell$, and the regular formulas identify $b,\mu,\nu,q$;
$s,V_i$ follow. The scalar covariance inverse also separates zero scale.
For a contrast estimate, $(c_1x+c_2y)^2\le2c_1^2x^2+2c_2^2y^2$ gives
upper order $c_1^2r_1+c_2^2r_2$.

The confidence set \eqref{eq:confidence-ball}, constructed over
$\mathcal D(s_0)$, has uniform coverage by the same moment MSE,
including at zero scale. Every candidate has moment error at most $Z_n$
about the truth off or on the coverage event. Equation
\eqref{eq:directional-inverse} bounds the coordinate diameters separately,
with each capped by its compact interval. The expected squared contrast diameter
is at most $C_\alpha(c_1^2r_1+c_2^2r_2)$. This derives the directional
confidence upper claim without a plug-in leverage estimate.

\subsection{All-price lower bound for the exposed persistence}

For $a_1$, we restrict to $q=0$, $v_1=v_2=s$, fixed signed $\ell$, fixed
$a_2$ and fixed admissible regular nuisance parameters. With $\varrho=1-\ell^2$,
we use the representation
\[
 B_1=\ell U_2+\sqrt{\varrho}W_1
\]
with $U_1,U_2,W_1,B_2$ independent standard Brownian motions. The enlarged
experiment observes the entire prehistory of $U_2$ and the $Y_2$ path,
together with every clean asset-1 increment of the full grid. The law of
the two paths is independent of $a=a_1$. The asset-2 clean levels can be reconstructed from
$U_2,Y_2$ with their fixed starting level; adding the original level
errors to both assets is a parameter-independent kernel. The target
$Y_1$ stays hidden.

Given $U_2$, its conditional mean and covariance are
\begin{equation}\label{eq:a1-conditional}
 m_a(t)=\mu_1+\ell\sqrt{2as}\int_{-\infty}^t e^{-a(t-u)}dU_2(u),
 \qquad s\varrho e^{-a|t-u|}.
\end{equation}
Since $U_1$ is independent of this state, the emissions conditional on
$Y_1$ have fixed means $b_1|I_j|$ and variances
$D_j=\int_{I_j}e^{2Y_1(t)}dt$.

For finite quadratures put
\[
 J_j^{\mathrm m}=\sum_{u\in j}w_u\dot m_a(t_u)e^{2y_u},\qquad
 W_{jk}=\sum_{u\in j,v\in k}w_uw_v\dot R_a(t_u,t_v)e^{2y_u+2y_v},
 \quad R_a(t,u)=e^{-a|t-u|}.
\]
With $S_j=H_2(Z_j)/(2D_j)$ and
$T_j^{\mathrm H}=H_4(Z_j)/(4D_j^2)$, define
\[
 \mathcal Z=\sum_jW_{jj}T_j^{\mathrm H}
                  +2\sum_{j<k}W_{jk}S_jS_k.
\]
Equation \eqref{eq:Gaussian-diff} gives the complete derivative
\begin{equation}\label{eq:a1-score}
 \partial_a p=\EE\left[k\left\{2\sum_jJ_j^{\mathrm m}S_j
                              +2s\varrho\mathcal Z\right\}\mid U_2\right].
\end{equation}
Here $\dot R_a(t,t)=0$, so there is no covariance-diagonal derivative
term in this scalar fixed-variance submodel. The state-mean derivative
term is retained. Under $k$, its singleton $H_2$ supports are orthogonal
to the covariance singleton $H_4$ and pair $H_2H_2$ supports. Since
$|W_{jk}|\le a_{jk}D_jD_k$, where
$a_{jk}=\sup_{I_j,I_k}|t-u|e^{-a|t-u|}$, Jensen and Hermite norms give
\begin{equation}\label{eq:a1-information}
 I(a)\le2\sum_j\EE[(J_j^{\mathrm m})^2/D_j^2]
                         +Cs^2\varrho^2\sum_{j,k}a_{jk}^2.
\end{equation}
For instance $\EE_k(T_j^{\mathrm H})^2=3/(2D_j^4)$ and
$\EE_k S_j^2=1/(2D_j^2)$, which cancel the $D$ factors in the covariance
part.

Differentiating the mean convolution gives
\begin{equation}\label{eq:a1-mean-derivative}
 \dot m_a(t)=\ell\sqrt{2as}\int_{-\infty}^t
       [1/(2a)-(t-u)]e^{-a(t-u)}dU_2(u),
 \qquad \EE\dot m_a(t)^2=\frac{s\ell^2}{4a^2}.
\end{equation}
The variance identity follows by integrating
$2as\ell^2\int_0^\infty(1/(2a)-v)^2e^{-2av}dv$.
For the bounded-interval supremum required to control $J_j^{\mathrm m}$,
let $x=m_a-\mu_1$. It solves the stable two-equation system
\[
 dx=-ax\,dt+\ell\sqrt{2as}\,dU_2,
 \qquad d\dot m_a=-(x+a\dot m_a)dt
                         +\ell\sqrt{s/(2a)}\,dU_2.
\]
Its stationary starting second moments are $\EE x^2=s\ell^2$ and those
in \eqref{eq:a1-mean-derivative}. On an interval of length at most $d_+$,
the integrated equations, Cauchy--Schwarz for drift integrals and Doob's
$L^2$ inequality for the Brownian terms give an inequality of the form
\[
 \EE\sup_{t_0\le t\le t_0+u}(|x(t)|^2+|\dot m_a(t)|^2)
 \le Cs\ell^2+C\int_0^u
       \EE\sup_{t_0\le t\le t_0+v}(|x(t)|^2+|\dot m_a(t)|^2)dv.
\]
Gronwall yields $\EE\sup_{I_j}|\dot m_a|^2\le Cs\ell^2$, uniformly on
the positive compact $a$ interval. Positive weights imply
$|J_j^{\mathrm m}|/D_j\le\sup_{I_j}|\dot m_a|$.
The same geometric summability as \eqref{eq:summability} bounds
$\sum a_{jk}^2$ by $Cn$. Thus \eqref{eq:a1-information} gives
$I(a)\le Cn(s^2+s\ell^2)$ uniformly in quadrature. The parameter
integration and exact-kernel convergence proved after
\eqref{eq:finite-information} apply here, including the additional observed paths.
After reconstruction,
\begin{equation}\label{eq:a1-Hellinger}
 \Hell^2(P_a^{(n)},P_{a'}^{(n)})
                  \le Cn(s^2+s\ell^2)|a-a'|^2.
\end{equation}
An interior pair with separation $c\sqrt{r_1}$ therefore proves the
all-price $r_1$ lower order.

\subsection{All-price lower bound for the unexposed persistence}

We now fix $a_1$ and vary $a=a_2$, again with $q=0$, $v_i=s$, fixed signed
$\ell$ and nuisance parameters. Here the enlarged experiment observes the
parameter-free $B_1,U_1$ on their full prehistory, together with all clean
asset-2 increments. These observations reconstruct asset-1 prices using the
fixed $a_1$. We use the representation
\[
 U_2=\ell B_1+\sqrt{\varrho}W_2
\]
with $B_1,W_2,B_2,U_1$ independent. The hidden $Y_2$ is independent of
$B_1$, and the exact conditional emission means and variances are
\begin{equation}\label{eq:a2-emission}
 a_j=b_2|I_j|+\ell\int_{I_j}e^{Y_2(t)}dB_1(t),\qquad
 d_j=\varrho D_j,\quad D_j=\int_{I_j}e^{2Y_2(t)}dt.
\end{equation}
The state dependence of the emission means requires retaining their derivatives.

In positive left quadratures, the $B_1$ increments $z_u$ are independent,
with variance $w_u$. The emission derivatives are now
\[
 \mathfrak m_u=\ell e^{y_u}z_u,\qquad
 \mathfrak d_u=2\varrho w_ue^{2y_u}.
\]
The state mean is fixed and its covariance is $sR_a$,
$R_a(t,u)=e^{-a|t-u|}$. Only in this scalar fixed-variance model is
$\dot R_a(t,t)=0$. All four arrays
$\mathsf C^{mm},\mathsf C^{md},\mathsf C^{dm},\mathsf C^{dd}$ from
\eqref{eq:score-pairs} enter with these derivatives. The exact marginal
score derivative is
\begin{equation}\label{eq:a2-score}
 \partial_a p=(s/2)\EE[k\mathcal Q\mid B_1],
\end{equation}
where $\mathcal Q$ is the full expression \eqref{eq:Q-score}; there are
no mean-state or covariance-diagonal additions in this particular
submodel. All singleton degrees and pair supports in $\mathcal Q$ are
orthogonal here. The Hermite identities give the exact norm
\begin{align}\label{eq:a2-score-norm}
 \int k\mathcal Q^2\,dx={}&\sum_j\left[
 \frac{2(\mathsf C^{mm}_{jj})^2}{d_j^2}
 +\frac{6(\mathsf C^{md}_{jj})^2}{d_j^3}
 +\frac{3(\mathsf C^{dd}_{jj})^2}{2d_j^4}\right]\notag\\
 &+4\sum_{j<k}\left[
 \frac{(\mathsf C^{mm}_{jk})^2}{d_jd_k}
 +\frac{(\mathsf C^{md}_{jk})^2}{2d_jd_k^2}
 +\frac{(\mathsf C^{dm}_{jk})^2}{2d_j^2d_k}
 +\frac{(\mathsf C^{dd}_{jk})^2}{4d_j^2d_k^2}\right].
\end{align}
Here the $mm$, $md$ and $dm$ terms account for the random mean, and the
$dd$ terms for the variance mixture.

The integrated-variance denominators cancel as follows. Let
\[
 a_{jk}=\sup_{I_j,I_k}|\dot R_a(t,u)|.
\]
Conditional on the full
$Y_2$ path, the $z_u$ are independent centered Gaussians, since
$Y_2$ is independent of $B_1$ in this richer experiment. Their second and
fourth moments give
\begin{equation}\label{eq:a2-coefficient-bounds}
 \begin{aligned}
 \EE_{B_1}(\mathsf C^{mm}_{jj})^2
       &\le2\ell^4a_{jj}^2D_j^2,\\
 \EE_{B_1}(\mathsf C^{mm}_{jk})^2
       &\le\ell^4a_{jk}^2D_jD_k\quad(j\ne k),\\
 \EE_{B_1}(\mathsf C^{md}_{jk})^2
       &\le4\ell^2\varrho^2 a_{jk}^2D_jD_k^2,\\
 |\mathsf C^{dd}_{jk}|&\le4\varrho^2 a_{jk}D_jD_k.
 \end{aligned}
\end{equation}
The transposed $dm$ inequality interchanges $j,k$. To see the first
bound, $\mathsf C^{mm}_{jj}$ is a Gaussian quadratic form with zero
conditional mean, since its diagonal coefficients $\dot R_a(t_u,t_u)$
vanish. Its variance is twice the weighted sum of squared off-diagonal
coefficients:
\[
 2\ell^4\sum_{u,v\in j}\dot R_a(t_u,t_v)^2
                 e^{2y_u+2y_v}w_uw_v\le2\ell^4a_{jj}^2D_j^2.
\]
For distinct emissions the two independent Brownian sums give the second
bound without the factor two. For $md$, conditional variance of the
linear Gaussian sum is
\[
 \ell^2\sum_{u\in j}w_ue^{2y_u}
    \left(2\varrho\sum_{v\in k}\dot R_a(t_u,t_v)w_ve^{2y_v}\right)^2
 \le4\ell^2\varrho^2a_{jk}^2D_jD_k^2.
\]
The $dd$ bound is the absolute bound for two deterministic weighted sums.
These bounds hold uniformly in the number of quadrature nodes within each emission.

After substitution of $d_j=\varrho D_j$ into \eqref{eq:a2-score-norm}, the three
singleton contributions are bounded respectively by
$4\ell^4/\varrho^2$, $24\ell^2/\varrho$, and $24$, times $a_{jj}^2$.
The analogous pair bounds are controlled by the same fixed envelope.
Consequently
\begin{equation}\label{eq:a2-cancellation}
 \EE\int k\mathcal Q^2\,dx
 \le24\left(1+\frac{\ell^2}{\varrho}
                +\frac{\ell^4}{\varrho^2}\right)\sum_{j,k}a_{jk}^2
 \le\frac{24}{\varrho^2}\sum_{j,k}a_{jk}^2.
\end{equation}
The last inequality follows from
$\varrho^2+\ell^2\varrho+\ell^4\le1$ with $\varrho=1-\ell^2$.
There are no negative powers of $D_j$ left. Jensen for
\eqref{eq:a2-score}, summability and $\varrho\ge\kappa$ give
\begin{equation}\label{eq:a2-information}
 I(a)\le Cns^2/(1-\ell^2)^2\le C_\kappa ns^2.
\end{equation}
Finite quadratures, parameter integration and convergence to exact
conditional kernels are justified exactly as in the proof of
Lemma~\ref{lem:all-price}. Here the target state is independent of the
additional observed driver, but the emission mean remains stochastic. Reconstruction
of asset-1 prices and addition of one error at each distinct time to both
assets give
\begin{equation}\label{eq:a2-Hellinger}
 \Hell^2(P_a^{(n)},P_{a'}^{(n)})\le C_\kappa ns^2|a-a'|^2.
\end{equation}
An interior $a_2$ pair with separation $c\sqrt{r_2}$ proves its all-price
lower order. In particular, fixed nonzero cross leverage does not
accelerate this persistence.

For $c_1a_1+c_2a_2$, each coordinate can vary independently over an interval
of positive length. The two testing arguments lower-bound worst-case risk
by constants times $c_1^2r_1$ and $c_2^2r_2$; their maximum is at least
half their sum. The upper bound was proved by the same coordinatewise
argument after \eqref{eq:directional-inverse}. For honest sets,
\eqref{eq:two-point-diameter} applies separately to these same two-point
submodels. The maximum-to-sum inequality then gives, for fixed
$\alpha<1/4$, the matching expected squared diameter lower
order. All infima here permit every estimator using the full noisy grid;
fixing nuisance parameters on a lower submodel does not restrict the procedure.
This completes Theorem~\ref{thm:diagonal}. None of these rate bounds is
extended to perfect leverage; the constant in
\eqref{eq:a2-information} explicitly uses its positive residual margin.

\section{Conclusion}

The precision warranted by noisy prices depends on where leverage carries
information through the hidden dynamics. In the coupled rank-one example,
three exact lags identify all four drift entries. In the separate diagonal
known-support example, the second persistence remains limited by
covariance-only information. The rates quantify this difference even when
all nuisance parameters are unknown. Their lower bounds use every noisy
price level, and honest confidence sets express the resulting uncertainty
without assuming excitation for coverage. The precision guarantees retain
fixed excitation and residual margins and a sufficiently small,
unevaluated covariance ceiling; they do not yet give calibrated sample
sizes or practical confidence widths.

\section*{Acknowledgements}

The author used ChatGPT (OpenAI, GPT-6 Astra) to assist in drafting and
polishing the exposition from the author's notes and proof sketches, as well
as in checking for possible gaps and inconsistencies. The author is
responsible for the final statements and proofs.

\begin{thebibliography}{9}
\enlargethispage{2\baselineskip}

\bibitem{ChengRenaultSangrey2024}
X.~Cheng, E.~Renault and P.~Sangrey,
\emph{Identifying the Volatility Risk Price Through the Leverage Effect},
author manuscript, 23 April 2024.
Publication DOI: \texttt{10.1016/j.jeconom.2024.105943}.

\bibitem{AhsanDufour2024}
Md.~N.~Ahsan and J.-M.~Dufour,
\emph{High-frequency instruments and identification-robust inference for
stochastic volatility models}, October 2024 author manuscript, including
Supplementary Appendix.
Journal publication (2025) DOI: \texttt{10.1111/jtsa.12812}.

\bibitem{ADR2025journal}
Md.~N.~Ahsan, J.-M.~Dufour and G.~Rodriguez-Rondon,
\emph{Estimation and Inference for Higher-Order Stochastic Volatility
Models With Leverage}, Journal of Time Series Analysis \textbf{46}(6)
(2025), 1064--1084.
DOI: \texttt{10.1111/jtsa.12851}.

\bibitem{GregoirMeddahi2025}
St\'ephane Gregoir and N.~Meddahi,
\emph{Towards Identification of Shocks in Linear State-Space Models:
Application to Stochastic Volatility Model}, Journal of Time Series
Analysis \textbf{46}(6) (2025), 1205--1223; first published 8 September
2025. DOI: \texttt{10.1111/jtsa.70012}.

\bibitem{OymakOzay2019}
S.~Oymak and N.~Ozay,
\emph{Non-asymptotic Identification of LTI Systems from a Single
Trajectory}, arXiv:1806.05722v2, 3 February 2019.

\bibitem{MasoliverPerello2005}
J.~Masoliver and J.~Perell\'o,
\emph{Multiple time scales and the exponential Ornstein--Uhlenbeck
stochastic volatility model}, arXiv:cond-mat/0501639v1, 26 January 2005;
Quantitative Finance \textbf{6} (2006), 423--433.
DOI: \texttt{10.1080/14697680600727547}.

\bibitem{Gloter2007}
A.~Gloter,
\emph{Efficient estimation of drift parameters in stochastic volatility
models}, Finance and Stochastics \textbf{11} (2007), 495--519.
DOI: \texttt{10.1007/s00780-007-0048-2}.

\bibitem{AbrahamNauletGassiat2022}
K.~Abraham, Z.~Naulet and \'E.~Gassiat,
\emph{Fundamental limits for learning hidden Markov model parameters},
arXiv:2106.12936v3, 24 October 2022.

\bibitem{ChongEtAl2024}
C.~H.~Chong, M.~Hoffmann, Y.~Liu, M.~Rosenbaum and G.~Szymanski,
\emph{Statistical inference for rough volatility: Minimax theory},
arXiv:2210.01214v2, 15 February 2024.

\end{thebibliography}
\end{document}
